%% toc: P3 37 37.1
\section{Proofs for abstract logics, principal examples, and comparison maps}
\label{sec:annex-p03-w03-p03-047-abstract-logics}

Throughout, structures have nonempty domains and the ambient set theory is
\(\mathsf{ZFC}\).  For a logic in the sense of
\cref{def:p03-w03-p03-047-ebbinghaus-logic}, a consequence
\(\varphi_0\models\varphi_1\) between \(\varphi_i\in\mathscr L[\tau_i]\) is
evaluated over \(\operatorname{Str}[\tau_0\cup\tau_1]\); by the reduct
property the outcome is the same over any larger vocabulary, since every
structure for a larger vocabulary satisfies each \(\varphi_i\) exactly when
its reduct does, and every \((\tau_0\cup\tau_1)\)-structure has an expansion
to the larger vocabulary (interpret the new symbols arbitrarily on the
nonempty domains).  The same remark applies to negations and conjunctions
supplied by the closure properties: if \(\neg\varphi\in\mathscr L[\tau]\) and
\(\tau\subseteq\tau_{\mathfrak B}\), then \(\mathfrak B\models\neg\varphi\)
iff \(\mathfrak B\restriction\tau\models\neg\varphi\) iff
\(\mathfrak B\restriction\tau\not\models\varphi\) iff
\(\mathfrak B\not\models\varphi\).

\begin{annexitem}{annex:p03-w03-p03-047-abstract-completeness}{thm:p03-w03-p03-047-abstract-completeness}{Validity and effective calculi}
\begin{proof}
(1) Let \(\mathsf{Ax}\) be the recursively enumerable set of axioms and
\(\mathsf{Ru}\) the recursively enumerable set of rule instances, each a pair
consisting of a finite set of premisses and a conclusion.  A finite sequence
\(\langle\chi_1,\ldots,\chi_m\rangle\) is a derivation iff for each
\(i\le m\) either \(\chi_i\in\mathsf{Ax}\) or some instance in \(\mathsf{Ru}\)
has conclusion \(\chi_i\) and all premisses among
\(\chi_1,\ldots,\chi_{i-1}\).  Each of the \(m\) conditions is the
existence of a stage in the enumeration of \(\mathsf{Ax}\) or of
\(\mathsf{Ru}\) together with a decidable check, so the set of derivations is
recursively enumerable: run the two enumerations and accept a sequence once
all \(m\) conditions have been witnessed.  The set of last members of
derivations is the image of this set under a computable map, hence
recursively enumerable.  By hypothesis this image is the set of valid
sentences.

(2) Let \(V\) be the set of valid sentences, recursively enumerable by
hypothesis, and take \(\mathsf{Ax}=V\) and \(\mathsf{Ru}=\varnothing\).  This
is an effective calculus.  In it every member of a derivation is an axiom,
so the derivable sentences are exactly the members of \(V\): each
\(\chi\in V\) is derived by the one-term sequence \(\langle\chi\rangle\), and
nothing outside \(V\) occurs in a derivation.  Every axiom is valid, so the
calculus is sound, and it derives every valid sentence.
\end{proof}
\end{annexitem}

\begin{annexitem}{annex:p03-w03-p03-047-completeness-compactness}{thm:p03-w03-p03-047-completeness-compactness}{From completeness and finite derivations to compactness}
\begin{proof}
Let \(T\) be a set of sentences without a model.  Then, vacuously, every
model of \(T\) is a model of \(\bot\), so strong completeness gives
\(T\vdash\bot\).  By finite character some finite \(T_0\subseteq T\) has
\(T_0\vdash\bot\).  By soundness every model of \(T_0\) is a model of
\(\bot\); since \(\bot\) has no model, \(T_0\) has no model.  Contraposing,
if every finite subset of \(T\) has a model then \(T\) has one: this is the
strong compactness property.  If completeness is assumed only for countable
sets, apply the same argument to a countable \(T\); the finite \(T_0\) is
then a finite subset of a countable set, and the conclusion is the countable
compactness property.
\end{proof}
\end{annexitem}

\begin{annexitem}{annex:p03-w03-p03-047-fo-undefinable}{thm:p03-w03-p03-047-fo-undefinable}{Undefinability of finiteness and uncountability}
\begin{proof}
For \(n\ge1\) let \(\sigma_n\) be the first-order sentence
\(\exists x_1\ldots x_n\bigl(\bigwedge_{i<j\le n}x_i\ne x_j\land
\bigwedge_{i\le n}Px_i\bigr)\), true exactly when \(P\) has at least \(n\)
elements.

(1) Suppose \(\psi\) is first-order and \(\mathfrak A\models\psi\) iff
\(P^{\mathfrak A}\) is infinite.  Let
\(T=\{\neg\psi\}\cup\{\sigma_n:n\ge1\}\), a countable set.  A finite subset
of \(T\) is contained in \(\{\neg\psi,\sigma_1,\ldots,\sigma_m\}\) for some
\(m\); the structure with domain \(\{0,\ldots,m-1\}\) and \(P\) the whole
domain satisfies \(\sigma_1,\ldots,\sigma_m\), and since its \(P\) is finite
it satisfies \(\neg\psi\).  By first-order compactness
(\cref{thm:p02-w02-p02-005-bw-compactness}) applied to the countable set \(T\), some
\(\mathfrak A\models T\).  Since \(\mathfrak A\models\sigma_n\) for all
\(n\), \(P^{\mathfrak A}\) is infinite, so \(\mathfrak A\models\psi\),
contradicting \(\mathfrak A\models\neg\psi\).  If instead \(\psi\) defined
finiteness of \(P\), then \(\neg\psi\) would define infiniteness, which was
just excluded.

(2) Suppose \(\psi\) is first-order and \(\mathfrak A\models\psi\) iff
\(P^{\mathfrak A}\) is uncountable.  The structure \(\mathfrak M\) with
domain \(\omega_1\) and \(P^{\mathfrak M}=\omega_1\) satisfies \(\psi\).  So
\(\{\psi\}\), a set of at most \(\aleph_0\) sentences, has a model, and by
the L\"owenheim--Skolem theorem (\cref{thm:p02-w02-p02-005-bw-ls}) it has a
model \(\mathfrak N\) whose domain has power at most \(\aleph_0\).  Then
\(\mathfrak N\models\psi\), yet \(P^{\mathfrak N}\subseteq N\) is countable,
a contradiction.
\end{proof}
\end{annexitem}

\begin{annexitem}{annex:p03-w03-p03-047-q0-omega}{thm:p03-w03-p03-047-q0-omega}{The \(\mathscr L(Q_0)\)-characterization of \((\omega,<)\)}
\begin{proof}
In \((\omega,<)\) the relation is a strict linear order, the domain is
infinite, and each \(n\) has exactly the \(n\) predecessors
\(0,\ldots,n-1\); so \(\theta\) holds, and by the isomorphism property it
holds in every copy of \((\omega,<)\).

Conversely let \((A,<)\models\theta\).  For \(a\in A\) put
\(\operatorname{pr}(a)=\{x\in A:x<a\}\), which is finite by the third
conjunct, and \(h(a)=|\operatorname{pr}(a)|\in\omega\).  If \(a<b\) then
\(\operatorname{pr}(a)\subseteq\operatorname{pr}(b)\) by transitivity and
\(a\in\operatorname{pr}(b)\setminus\operatorname{pr}(a)\) by irreflexivity,
so \(h(a)<h(b)\).  By linearity, distinct \(a,b\) are comparable, so \(h\) is
injective, and \(h(a)<h(b)\) implies \(a<b\) (otherwise \(b<a\) or \(b=a\)
would give \(h(b)\le h(a)\)).  The image of \(h\) is downward closed: if
\(h(b)=n\) list \(\operatorname{pr}(b)\) increasingly as
\(b_0<\cdots<b_{n-1}\); for \(m<n\), the predecessors of \(b_m\) are below
\(b\), hence in \(\operatorname{pr}(b)\), and are exactly
\(b_0,\ldots,b_{m-1}\), so \(h(b_m)=m\).  The second conjunct makes \(A\)
infinite, so the injective \(h\) has infinite downward closed image, which is
\(\omega\).  Thus \(h\) is an isomorphism of \((A,<)\) onto \((\omega,<)\).
\end{proof}
\end{annexitem}

\begin{annexitem}{annex:p03-w03-p03-047-q0-not-re}{thm:p03-w03-p03-047-q0-not-re}{No completeness theorem for \(\mathscr L(Q_0)\)}
\begin{proof}
(1) In \((\mathcal N,<)\) the order is the usual one, so \(\theta\) holds by
\cref{thm:p03-w03-p03-047-q0-omega}; nothing lies below \(0\); \(n+1\) is the
immediate successor of \(n\); and the four recursion equations are the
defining equations of addition and multiplication.  By the isomorphism
property every isomorphic copy satisfies \(\chi_{\mathcal N}\).

Conversely let \(\mathfrak A\models\chi_{\mathcal N}\).  By
\cref{thm:p03-w03-p03-047-q0-omega} there is an isomorphism
\(h:(A,<^{\mathfrak A})\to(\omega,<)\).  Since no element is below
\(0^{\mathfrak A}\) and the order is linear, \(0^{\mathfrak A}\) is the least
element, so \(h(0^{\mathfrak A})=0\).  For \(a\in A\), \(S^{\mathfrak A}a\) is
the immediate successor of \(a\) by the second displayed conjunct; \(h\)
carries immediate successors to immediate successors, and in \(\omega\) the
immediate successor of \(n\) is \(n+1\); so \(h(S^{\mathfrak A}a)=h(a)+1\).
Fix \(a\in A\) and prove \(h(a+^{\mathfrak A}b)=h(a)+h(b)\) by induction on
\(n=h(b)\).  If \(n=0\) then \(b=0^{\mathfrak A}\) and
\(a+^{\mathfrak A}0^{\mathfrak A}=a\).  If \(h(b)=n+1\), let
\(b'=h^{-1}(n)\); then \(h(S^{\mathfrak A}b')=n+1\), so
\(b=S^{\mathfrak A}b'\), and
\(h(a+^{\mathfrak A}b)=h(S^{\mathfrak A}(a+^{\mathfrak A}b'))=h(a+^{\mathfrak A}b')+1=h(a)+n+1\).
The same induction with the multiplication equations, using the result for
addition in the successor step
\(h(a\times^{\mathfrak A}S^{\mathfrak A}b')=h(a\times^{\mathfrak A}b')+h(a)\),
gives \(h(a\times^{\mathfrak A}b)=h(a)\cdot h(b)\).  So \(h\) is an
isomorphism of \(\mathfrak A\) onto \((\mathcal N,<)\).

(2) If \(\mathcal N\models\sigma\) and \(\mathfrak A\models\chi_{\mathcal N}\),
then \(\mathfrak A\restriction\Lang_A\cong\mathcal N\) by part~(1), so
\(\mathfrak A\models\sigma\); hence \(\chi_{\mathcal N}\to\sigma\) is valid.
If \(\chi_{\mathcal N}\to\sigma\) is valid, it holds in
\((\mathcal N,<)\), which satisfies \(\chi_{\mathcal N}\), so
\((\mathcal N,<)\models\sigma\) and \(\mathcal N\models\sigma\) since \(<\)
does not occur in \(\sigma\).

(3) Let \(r\) be the computable map of the hypothesis and \(V\) the set of
codes of valid sentences.  By part~(2), \(\mathrm{TA}=\{c:c\text{ codes a
first-order }\Lang_A\text{-sentence and }r(c)\in V\}\).  The set of codes of
first-order \(\Lang_A\)-sentences is decidable (clause~1 of
\cref{thm:p02-w02-p02-040-no-effective-calculus}).  If \(V\) were
recursively enumerable, then so would be \(\mathrm{TA}\): enumerate \(V\),
and accept a code \(c\) of a sentence once \(r(c)\) appears.  This
contradicts clause~3 of \cref{thm:p02-w02-p02-040-no-effective-calculus}.
So \(V\) is not recursively enumerable, and by part~1 of
\cref{thm:p03-w03-p03-047-abstract-completeness} no effective calculus
derives exactly the members of \(V\).
\end{proof}
\end{annexitem}

\begin{annexitem}{annex:p03-w03-p03-047-q1-mostowski-facts}{thm:p03-w03-p03-047-q1-mostowski-facts}{Expressive facts for \(\mathscr L(Q_1)\)}
\begin{proof}
Throughout, Keisler's standard models are the structures in which \(Q\) is
read as \(Q_1\) of \cref{def:p03-q-one}.

(1) The calculus of \cref{def:p02-w02c-p02-003-kq-calculus} has a decidable
set of axioms, by hypothesis, and the two rules modus ponens and
generalization, each with finitely many premisses and a decidable instance
relation; it is therefore an effective calculus in the sense of
\cref{def:p03-w03-p03-047-abstract-completeness}.  By
\cref{thm:p02-w02c-p02-003-keisler-completeness} a sentence is provable in
it iff it holds in every standard model, that is, iff it is valid in
\(\mathscr L(Q_1)\).  Part~1 of
\cref{thm:p03-w03-p03-047-abstract-completeness} gives recursive
enumerability.

(2) Suppose \(\psi\) were such a sentence.  Add a constant \(c\) and, for
\(n\ge1\), let \(\delta_n\) be the first-order sentence
\(\exists x_1\ldots\exists x_n(\bigwedge_{i<j}x_i\ne x_j\land\bigwedge_i x_i<c)\).
Let \(\Sigma=\{\psi\}\cup\{\delta_n:n\ge1\}\), a set of sentences of the
countable vocabulary \(\{<,c\}\).  A finite subset mentions only
\(\delta_n\) with \(n\le N\) for some \(N\); \((\omega,<,N)\) satisfies
\(\psi\) (it is isomorphic to \((\omega,<)\)) and has the \(N\) elements
\(0,\ldots,N-1\) below \(N\), so satisfies the subset.  By
\cref{thm:p02-w02c-p02-003-keisler-compactness} \(\Sigma\) has a model
\((A,<,c)\).  Then \((A,<)\models\psi\), so \((A,<)\cong(\omega,<)\) and
\(c\) has finitely many predecessors, say \(m\); but
\((A,<,c)\models\delta_{m+1}\), a contradiction.  If \(\mathscr L(Q_1)\)
included \(\mathscr L(Q_0)\), the sentence \(\theta\) of
\cref{thm:p03-w03-p03-047-q0-omega} would have an equivalent \(\psi\) of
this kind.

(3) Let \(\mathfrak G\) be a group.  If \(\mathfrak G\) is generated by a
countable set \(X\), every element is a finite product
\(x_1^{\epsilon_1}\cdots x_n^{\epsilon_n}\) with \(x_i\in X\) and
\(\epsilon_i\in\{1,-1\}\).  For each \(n\) the set of such words is the image
of the countable set \((X\times\{1,-1\})^n\), and \(G\) is the union over
\(n\in\omega\) of these countable images, hence countable (a countable union
of countable sets, by the axiom of choice).  So \(\mathfrak G\models\neg
Q_1x\,(x=x)\).  Conversely, if \(G\) is countable, \(G\) itself is a
countable generating set.  The clause for \(Q_1x\,(x=x)\) is the semantic
rule of \cref{def:p03-q-one} with \(\varphi\) the formula \(x=x\).
\end{proof}
\end{annexitem}

\begin{annexitem}{annex:p03-w03-p03-047-q0-vector-spaces}{thm:p03-w03-p03-047-q0-vector-spaces}{Finite dimension and \(Q_0\)}
\begin{proof}
(1) Let \(F\) have \(q\) elements.  A vector space \(V\) over \(F\) of finite
dimension \(n\) is isomorphic to \(F^n\) and has \(q^n\) elements.  If \(V\)
has infinite dimension it contains, for each \(n\), \(n\) linearly
independent vectors, whose span has \(q^n\) elements, so \(V\) is infinite.
Hence \(V\) is finite-dimensional iff it is finite iff
\(V\models\neg Q_0x\,(x=x)\).

(2) \emph{Step 1: linear terms.}  Every model \(V\) of \(T\) is infinite:
for \(v\ne0\) the vectors \(\lambda_a v\) (\(a\in\mathbb Q\)) are pairwise
distinct, since \(\lambda_av=\lambda_bv\) gives \(\lambda_{a-b}v=0\) and
\(a-b\ne0\) would give \(v=\lambda_{(a-b)^{-1}}\lambda_{a-b}v=0\).  By
induction on terms, the vector-space axioms in \(T\) prove every
\(\tau_{\mathbb Q}\)-term \(t(x,\bar y)\) equal to a term
\(\lambda_a x+\sum_i\lambda_{a_i}y_i\) with \(a,a_i\in\mathbb Q\) (a variable
\(z\) is \(\lambda_1z\); \(0\) is \(\lambda_0x\); sums, differences and
scalar multiples of such terms are again of this form after collecting
coefficients).  Hence every atomic formula \(t_1=t_2\) is \(T\)-equivalent
to \(\lambda_a x=s(\bar y)\) with \(s\) a term in \(\bar y\) alone, and, when
\(a\ne0\), to \(x=\lambda_{a^{-1}}s(\bar y)\); when \(a=0\) it is the
formula \(0=s(\bar y)\), in which \(x\) does not occur.

\emph{Step 2: quantifier elimination.}  We show that for each
quantifier-free \(\varphi(x,\bar y)\) the formula \(\exists x\,\varphi\) is
\(T\)-equivalent to a quantifier-free formula.  Put \(\varphi\) in
disjunctive normal form; \(\exists x\) distributes over disjunction, so let
\(\varphi\) be a conjunction of literals, and move the literals without
\(x\) outside \(\exists x\).  By Step~1 the remaining literals are
equations \(x=s_i(\bar y)\) and inequations \(x\ne s_j(\bar y)\).  If an
equation \(x=s_1(\bar y)\) occurs, \(\exists x\,\varphi\) is equivalent to
the result of replacing \(x\) by \(s_1(\bar y)\) throughout.  Otherwise
\(\varphi\) is \(\bigwedge_{j\le k}x\ne s_j(\bar y)\), which excludes at most
\(k\) vectors; since every model of \(T\) is infinite, \(\exists x\,\varphi\)
holds for every assignment and is equivalent to \(0=0\).  By induction on
formulas, using this step for \(\exists\), every first-order formula of
\(\tau_{\mathbb Q}\) is \(T\)-equivalent to a quantifier-free one with the
same free variables.

\emph{Step 3: one-variable sets.}  Let \(\psi(x,\bar y)\) be quantifier-free
in disjunctive normal form with \(m\) disjuncts and \(k\) literals in all.
Fix \(V\models T\) and \(\bar b\in V\).  By Step~1, a disjunct that contains
an equation with nonzero coefficient of \(x\) is satisfied by at most one
vector; a disjunct that does not is satisfied either by no vector (if one of
its literals without \(x\) fails at \(\bar b\)) or by all vectors outside a
set of at most \(k\) elements.  So the set
\(\psi(V,\bar b)=\{v:V\models\psi[v,\bar b]\}\) either has at most \(m\)
elements or has a complement of at most \(k\) elements; since \(V\) is
infinite, \(\psi(V,\bar b)\) is infinite iff it has more than \(m\)
elements.  Thus \(Q_0x\,\psi\) and the first-order formula
\(\exists^{>m}x\,\psi\), that is
\(\exists x_0\cdots\exists x_m(\bigwedge_{i<j\le m}x_i\ne x_j\land\bigwedge_{i\le m}\psi(x_i,\bar y))\),
are satisfied by the same assignments in every model of \(T\).

\emph{Step 4: the translation.}  Define \(\chi^\circ\) by induction:
atomic formulas are unchanged, \(\circ\) commutes with \(\neg\), \(\land\)
and \(\exists\), and \((Q_0x\,\chi)^\circ=\exists^{>m}x\,\psi\), where
\(\psi\) is a quantifier-free formula in disjunctive normal form with \(m\)
disjuncts that is \(T\)-equivalent to \(\chi^\circ\) (Step~2).  Induction on
\(\chi\) shows that \(\chi\) and \(\chi^\circ\) are satisfied by the same
assignments in every model of \(T\): the atomic, Boolean and \(\exists\)
cases follow from the induction hypothesis, and in the \(Q_0\) case the set
of \(x\)-values satisfying \(\chi\) equals that satisfying \(\chi^\circ\),
hence that satisfying \(\psi\), and Step~3 applies.

(3) By Step~2 every first-order sentence of \(\tau_{\mathbb Q}\) is
\(T\)-equivalent to a quantifier-free sentence, that is, to a Boolean
combination of equations \(t_1=t_2\) between closed terms.  Every closed
term is built from \(0\) by \(+\), \(-\) and the \(\lambda_a\), and the
vector-space axioms prove it equal to \(0\); so each such equation holds in
every model of \(T\), and the truth value of the Boolean combination is the
same in all models.  Hence \(T\) is complete.  Both \(\mathbb Q\) and
\(\mathbb Q^{(\omega)}\) (finitely supported sequences of rationals) are
nonzero vector spaces over \(\mathbb Q\), so they are models of \(T\).  For an
\(\mathscr L(Q_0)\)-sentence \(\chi\), part~(2) gives
\(\mathbb Q\models\chi\) iff \(\mathbb Q\models\chi^\circ\) iff
\(\mathbb Q^{(\omega)}\models\chi^\circ\) (completeness) iff
\(\mathbb Q^{(\omega)}\models\chi\).  A sentence true in exactly the
finite-dimensional spaces would be true in \(\mathbb Q\) and false in
\(\mathbb Q^{(\omega)}\), which is impossible.
\end{proof}
\end{annexitem}

\begin{annexitem}{annex:p03-w03-p03-047-q0-beth-failure}{thm:p03-w03-p03-047-q0-beth-failure}{Implicit but not explicit definability in \(\mathscr L(Q_0)\)}
\begin{proof}
Write \(\mathrm{Ev}\subseteq\omega\) for the set of even numbers.

(1) Let \(\mathfrak A=(A,<,E)\models\varphi\).  By
\cref{thm:p03-w03-p03-047-q0-omega} there is an isomorphism
\(f:(\omega,<)\to(A,<)\); it is unique, since \(f(n)\) must be the element
of \(A\) with exactly \(n\) predecessors.  We show by induction on \(n\)
that \(f(n)\in E\) iff \(n\in\mathrm{Ev}\).  The element \(f(0)\) has no
predecessor, so the second conjunct puts it in \(E\).  For the step,
\(f(n+1)\) is the immediate successor of \(f(n)\) in \((A,<)\) because
\(f\) is an isomorphism and \(n+1\) is the immediate successor of \(n\) in
\(\omega\); the third conjunct gives \(f(n+1)\in E\) iff \(f(n)\notin E\)
iff \(n\notin\mathrm{Ev}\) iff \(n+1\in\mathrm{Ev}\).  Hence
\(E=f[\mathrm{Ev}]\) is determined by \((A,<)\), and two expansions of the
same \(\{<\}\)-structure to models of \(\varphi\) coincide.  Conversely
\((\omega,<,\mathrm{Ev})\models\varphi\): \(\theta\) holds by
\cref{thm:p03-w03-p03-047-q0-omega}, \(0\) is the only element without
predecessor and is even, and \(S(m,k)\) holds in \((\omega,<)\) exactly when
\(k=m+1\), which has the opposite parity.

(2) \emph{Step 1: elimination of \(Q_0\) over \((\omega,<)\).}  For an
\(\mathscr L(Q_0)\)-formula \(\chi\) of vocabulary \(\{<\}\cup C\), with
\(C\) a finite set of constants, define a first-order formula \(\chi^*\) with
the same free variables by leaving atomic formulas unchanged, commuting with
\(\neg\), \(\land\) and \(\exists\), and putting
\((Q_0x\,\chi)^*=\forall y\,\exists x\,(y<x\land\chi^*)\) with \(y\) a
variable not occurring in \(\chi\).  We show by induction on \(\chi\) that
for every expansion \(\mathfrak N\) of \((\omega,<)\) interpreting \(C\) and
every assignment \(s\) of the free variables,
\(\mathfrak N\models\chi[s]\) iff \(\mathfrak N\models\chi^*[s]\).  The
atomic and Boolean cases and the \(\exists\) case are immediate from the
induction hypothesis.  For \(Q_0x\,\chi\): the set
\(X=\{n:\mathfrak N\models\chi[s(x/n)]\}\) equals
\(\{n:\mathfrak N\models\chi^*[s(x/n)]\}\) by the induction hypothesis, and
a subset of \(\omega\) is infinite iff it has an element above every
\(m\in\omega\).  The left side says \(X\) is infinite, the right side that
\(X\) is unbounded, so they agree.

\emph{Step 2: the first-order consequence.}  Suppose
\(\psi(c)\in\mathscr L(Q_0)[\{<,c\}]\) satisfies
\(E^{\mathfrak A}=\{a:(\mathfrak A\restriction\{<\},a)\models\psi(c)\}\)
for every \(\mathfrak A\models\varphi\).  Applied to
\((\omega,<,\mathrm{Ev})\) and combined with Step~1 (with \(c\) replaced by a
free variable \(x\)), the first-order formula \(\psi^*(x)\) defines
\(\mathrm{Ev}\) in \((\omega,<)\).  Hence \((\omega,<)\) satisfies the
first-order sentence
\[
 \delta_1:\quad\forall x\,\forall y\,\bigl(S(x,y)\to(\psi^*(y)\leftrightarrow\neg\psi^*(x))\bigr),
\]
the first-order sentence \(\delta_2\) saying that \(<\) is a strict linear
order with a least element in which every element has an immediate
successor and every element other than the least has an immediate
predecessor,
and, for each \(n\in\omega\), the sentence \(\delta_{3,n}\) saying that every
element with at most \(n\) predecessors has exactly \(k\) predecessors for
some \(k\le n\), and that for each \(k\le n\) there is exactly one element
with exactly \(k\) predecessors.  All of these are first-order.

\emph{Step 3: a proper elementary extension and a shift.}  By
\cref{thm:p02-w02-p02-033-upward-ls-inclusion} (with \(\kappa=\aleph_0\))
there is a structure \(\mathfrak M=(M,<)\) with
\((\omega,<)\preccurlyeq\mathfrak M\) and \(M\ne\omega\).  It satisfies
\(\delta_1\), \(\delta_2\) and every \(\delta_{3,n}\).  Put
\(B=M\setminus\omega\ne\varnothing\).

(a) \(\omega\) is an initial segment of \(\mathfrak M\).  Let \(n\in\omega\)
and \(a\in M\) with \(a<n\).  The element \(n\) has exactly \(n\)
predecessors in \(\omega\), and since the formula ``\(x\) has exactly \(n\)
predecessors'' is first-order and \(\preccurlyeq\) preserves it with the
parameter \(n\), \(n\) has exactly \(n\) predecessors in \(\mathfrak M\),
namely \(0,\ldots,n-1\).  So \(a\in\omega\).

(b) For \(a\in B\) let \(a^+\) be its immediate successor in \(\mathfrak M\)
(by \(\delta_2\)).  Then \(a^+\in B\), for otherwise \(a<a^+\in\omega\)
gives \(a\in\omega\) by (a).  Every \(a\in B\) has an immediate predecessor
\(a^-\): \(a\) is not the least element, since by \(\delta_{3,0}\) the least
element of \(\mathfrak M\) is the unique element with no predecessor, which
is \(0\in\omega\) by elementarity.  And \(a^-\in B\): if \(a^-=n\in\omega\),
then \(n+1\) is the immediate successor of \(n\) in \(\omega\), hence in
\(\mathfrak M\) by elementarity (the formula \(S(x,y)\) with parameters
\(n,n+1\)), so \(a=n+1\in\omega\), a contradiction.  Since immediate
successors and predecessors are unique in a linear order, \(a\mapsto a^+\)
is a bijection of \(B\) onto \(B\) with inverse \(a\mapsto a^-\).

(c) Define \(\sigma:M\to M\) by \(\sigma(n)=n\) for \(n\in\omega\) and
\(\sigma(a)=a^+\) for \(a\in B\).  By (b) \(\sigma\) is a bijection.  It is
order preserving: on \(\omega\) it is the identity; if \(n\in\omega\) and
\(a\in B\) then \(n<a\) by (a) and linearity, and \(n<a<a^+\); if
\(a<b\) in \(B\) then \(a^+\le b<b^+\), since \(a^+\) is the least element
above \(a\).  A bijective order-preserving map between linear orders
reflects the order as well, so \(\sigma\) is an automorphism of
\(\mathfrak M\).

\emph{Step 4: contradiction.}  Fix \(b\in B\).  By the isomorphism property
of first-order logic, \(\sigma\) is an isomorphism
\((\mathfrak M,b)\cong(\mathfrak M,b^+)\), so
\(\mathfrak M\models\psi^*[b]\) iff \(\mathfrak M\models\psi^*[b^+]\).  But
\(S(b,b^+)\) holds in \(\mathfrak M\), and \(\delta_1\) gives
\(\mathfrak M\models\psi^*[b^+]\) iff \(\mathfrak M\not\models\psi^*[b]\).
The two equivalences are incompatible, so no such \(\psi\) exists.

(3) By (1) and (2) the Beth property of
\cref{def:p03-w03-p03-047-ebbinghaus-properties}(ix) fails for the relation
symbol \(E\).  The logic \(\mathscr L(Q_0)\) is one-sorted and has the atom,
negation and conjunction properties, since its formulas are closed under
\(\neg\) and \(\land\) and contain the atomic formulas.  If it had the
interpolation property, \cref{thm:p03-w03-p03-047-craig-beth} would make
\(E\) explicitly definable relative to \(\varphi\), contradicting (2).
\end{proof}
\end{annexitem}

\begin{annexitem}{annex:p03-w03-p03-047-lomega1-q0}{thm:p03-w03-p03-047-lomega1-q0}{The infinitary definition of \(Q_0\)}
\begin{proof}
Fix a structure and values of the free variables of \(\varphi\) other than
\(y\), and let \(S\) be the set of elements satisfying \(\varphi\).  If \(S\)
is infinite, then for every \(n\) and any \(x_1,\ldots,x_n\) the set
\(S\setminus\{x_1,\ldots,x_n\}\) is nonempty, and any of its members
witnesses \(y\); so every conjunct holds.  If \(S\) is finite with \(k\)
elements, consider the conjunct for \(n=k\) and let \(x_1,\ldots,x_k\)
enumerate \(S\) (for \(k=0\) the conjunct is \(\exists y\,\varphi(y)\)).  No
\(y\) satisfies \(\varphi\) while differing from all of them, so that
conjunct fails.  Hence the conjunction holds iff \(S\) is infinite, that is,
iff \(Q_0y\,\varphi(y)\) holds.

Translate \(\mathscr L(Q_0)\)-formulas by recursion: atomic formulas,
connectives and \(\exists,\forall\) are kept, and \(Q_0x\,\varphi\) goes to
the displayed conjunction for the translation of \(\varphi\) (with \(x\)
renamed to \(y\) and \(x_1,\ldots,x_n\) chosen outside the free variables).
The result has finitely many free variables and only countable
conjunctions, so it is an \(\mathscr L_{\omega_1\omega}\)-formula.  Induction
on formulas, using the first paragraph at the \(Q_0\) step, shows that each
formula and its translation are satisfied by the same structures and
assignments.
\end{proof}
\end{annexitem}

\begin{annexitem}{annex:p03-w03-p03-047-mostowski-general}{thm:p03-w03-p03-047-mostowski-general}{Self-defined satisfaction and the diagonal}
\begin{proof}
\emph{Preliminary: first-order sentences.}  If \(\mathscr L\) has the atom,
Boolean and particularization properties, every first-order sentence
\(\alpha\) of a vocabulary \(\tau\) has an \(\mathscr L\)-equivalent
\(\alpha^{\mathscr L}\in\mathscr L[\tau]\).  By induction on first-order
formulas \(\alpha(x_1,\dots,x_k)\) we show that for all new constants
\(d_1,\dots,d_k\) the sentence \(\alpha(d_1,\dots,d_k)\) of
\(\tau\cup\{d_1,\dots,d_k\}\) has an \(\mathscr L\)-equivalent: atomic
formulas by the atom property, \(\neg\) and \(\land\) by the Boolean
property, and \(\exists x\,\alpha(x,\bar d)\) by applying the
particularization property to the \(\mathscr L\)-equivalent of
\(\alpha(d,\bar d)\) and the constant \(d\).

(1) \(U\) is implicitly defined by \(\sigma\) by clause~(c) of
\cref{def:p03-w03-p03-047-self-satisfaction}.  Suppose that
\(\psi(c)\in\mathscr L[\tau_<\cup\{c\}]\) defines \(U\) explicitly relative
to \(\sigma\).  Since \((\mathcal N,<,U^\ast)\models\sigma\),
\[
 U^\ast=\{n\in\omega:(\mathcal N,<,n)\models\psi(c)\}.\tag{\(\dagger\)}
\]
Let \(d\) be a constant different from \(c\).  By the renaming property
(\cref{def:p03-w03-p03-047-ebbinghaus-logic}) there is
\(\psi(d)\in\mathscr L[\tau_<\cup\{d\}]\) with
\((\mathfrak A,a)\models\psi(d)\) iff \((\mathfrak A,a)\models\psi(c)\), the
constant being interpreted by \(a\) on both sides.  Take the
\(\mathscr L\)-equivalent of the first-order sentence \(\gamma(c,c,d)\) (the
preliminary), form with the Boolean property a sentence equivalent to
\(\gamma(c,c,d)\land\neg\psi(d)\), and apply the particularization property
to \(d\).  The result is \(\delta\in\mathscr L[\tau_<\cup\{c\}]\) with, for
every \(n\),
\[
 (\mathcal N,<,n)\models\delta
 \iff (\mathcal N,<,g(n,n))\not\models\psi(c),
\]
because \(\gamma\) defines the graph of \(g\) in \(\mathcal N\), so the only
witness for \(d\) is \(g(n,n)\).  By clause~(a), \(\delta\) has the same
models as \(\psi_e\) for some \(e\in\mathrm{Fm}\).  Then, using clause~(c)
and \((\dagger)\),
\[
 g(e,e)\in U^\ast\iff(\mathcal N,<,e)\models\psi_e
 \iff(\mathcal N,<,e)\models\delta
 \iff(\mathcal N,<,g(e,e))\not\models\psi(c)
 \iff g(e,e)\notin U^\ast ,
\]
a contradiction.  So \(U\) is not explicitly definable relative to
\(\sigma\), and \(\mathscr L\) lacks the Beth property.  If \(\mathscr L\)
had the interpolation property, \cref{thm:p03-w03-p03-047-craig-beth} (whose
hypotheses, the atom, negation and conjunction properties, hold here) would
make \(U\) explicitly definable; so interpolation fails as well.

(2) \emph{Coding.}  Let
\(\pi(x,y)=\tfrac12(x+y)(x+y+1)+y\), a bijection of \(\omega\times\omega\)
onto \(\omega\) with \(\pi(x,y)\ge x\) and \(\pi(x,y)\ge y\), whose graph is
defined in \(\mathcal N\) by the atomic formula
\(SS0\times z=((x+y)\times S(x+y))+(SS0\times y)\).  Code finite sequences
by \(\langle\rangle=0\) and
\(\langle a_0,a_1,\dots,a_k\rangle=\pi(a_0,\langle a_1,\dots,a_k\rangle)+1\);
then every entry of a sequence is smaller than its code, and the functions
\(s\mapsto(s)_i\) (the \(i\)-th entry, \(0\) if there is none) and
\((s,i,m)\mapsto s[i/m]\) (the code of the sequence with \(i\)-th entry
replaced by \(m\), the sequence being lengthened by zeros if needed) are
computable.  Formulas of \(\mathscr L(Q_0)[\tau_<\cup\{c\}]\) are built from
the variables \(v_0,v_1,\dots\) and are coded by
\(\ulcorner v_i\urcorner=\langle0,i\rangle\),
\(\ulcorner0\urcorner=\langle1\rangle\),
\(\ulcorner c\urcorner=\langle2\rangle\),
\(\ulcorner St\urcorner=\langle3,\ulcorner t\urcorner\rangle\),
\(\ulcorner t+t'\urcorner=\langle4,\ulcorner t\urcorner,\ulcorner t'\urcorner\rangle\),
\(\ulcorner t\times t'\urcorner=\langle5,\ulcorner t\urcorner,\ulcorner t'\urcorner\rangle\),
\(\ulcorner t=t'\urcorner=\langle6,\ulcorner t\urcorner,\ulcorner t'\urcorner\rangle\),
\(\ulcorner t<t'\urcorner=\langle7,\ulcorner t\urcorner,\ulcorner t'\urcorner\rangle\),
\(\ulcorner\neg\varphi\urcorner=\langle8,\ulcorner\varphi\urcorner\rangle\),
\(\ulcorner\varphi\land\varphi'\urcorner=\langle9,\ulcorner\varphi\urcorner,\ulcorner\varphi'\urcorner\rangle\),
\(\ulcorner\exists v_i\varphi\urcorner=\langle10,i,\ulcorner\varphi\urcorner\rangle\) and
\(\ulcorner Q_0v_i\varphi\urcorner=\langle11,i,\ulcorner\varphi\urcorner\rangle\).
For \(\mathscr L^{w2}\) add relation variables \(X^k_j\) (\(k\ge1\)), with
\(\ulcorner X^k_jt_1\cdots t_k\urcorner=\langle12,k,j,\langle\ulcorner t_1\urcorner,\dots,\ulcorner t_k\urcorner\rangle\rangle\)
and \(\ulcorner\exists X^k_j\varphi\urcorner=\langle13,k,j,\ulcorner\varphi\urcorner\rangle\),
in place of the clause for \(Q_0\).
Every component of a code is smaller than the code, so by course-of-values
recursion the following are decidable: the set of codes of formulas; the
relations \(\mathrm{Neg}(e,e')\), \(\mathrm{Conj}(e,e_1,e_2)\),
\(\mathrm{Ex}(e,i,e')\), \(\mathrm{Qz}(e,i,e')\) and
\(\mathrm{Ex2}(e,k,j,e')\) saying that \(e\) codes the negation, the
conjunction, the \(\exists v_i\)-, \(Q_0v_i\)- or \(\exists X^k_j\)-formula
of the formulas coded by \(e'\) or \(e_1,e_2\); and the set \(\mathrm{At}\)
of codes of atomic formulas.

An \emph{assignment code} is any \(s\in\omega\), read as follows: \(v_i\)
receives \(((s)_0)_i\), and \(X^k_j\) receives the finite relation coded by
\((((s)_1)_k)_j\), where a number \(r=\sum_{a\in A}2^a\) codes the finite
set \(A\subseteq\omega\) and, as a \(k\)-ary relation, the set of
\(k\)-tuples \((b_1,\dots,b_k)\) with \(\langle b_1,\dots,b_k\rangle\in A\).
Every assignment of numbers to finitely many individual variables and of
finite relations to finitely many relation variables is given by some
code.  The value of a term \(t\) under \(s\), with \(c\) read as \(n\), is a
computable function of \((\ulcorner t\urcorner,s,n)\); so the relation
\(\mathrm{True}_{\mathrm{at}}(e,s,n)\), ``\(e\) codes an atomic formula that
is true in \((\mathcal N,<)\) under \(s\) with \(c\) read as \(n\)'', is
decidable, and so are the relation \(\mathrm{Upd}(s,i,m,s')\), ``\(s'\) is
\(s\) with the entry for \(v_i\) replaced by \(m\)'', and the relation
\(\mathrm{Upd2}(s,k,j,r,s')\), ``\(s'\) is \(s\) with the entry for
\(X^k_j\) replaced by \(r\)''; each of the last two is the graph of a
function of its first arguments.  A decidable relation is in
\(\boldsymbol\Delta^0_1\) (\cref{thm:p04-ce-sigma-one}), so each of these
relations is defined in \(\mathcal N\) by a first-order \(\Lang_A\)-formula,
denoted by the same name.  Write \(U\langle e,s,n\rangle\) for
\(\exists w\,\exists z\,(w=\pi(s,n)\land z=\pi(e,w)\land Uz)\), the graph
of \(\pi\) being given by the atomic formula above.

\emph{Data.}  Let \(\mathscr L\) be \(\mathscr L(Q_0)\) or
\(\mathscr L^{w2}\).  Let \(\mathrm{Fm}\) be the set of codes of sentences of
\(\mathscr L[\tau_<\cup\{c\}]\), \(\psi_e\) the sentence coded by \(e\), and
\(g(e,n)=\pi(e,\pi(0,n))\), whose graph is defined by
\(\exists w\,(w=\pi(0,y)\land z=\pi(x,w))\).  Clause~(a) holds since every
sentence has a code.  Let \(U^\ast\) be the set of all \(\pi(e,\pi(s,n))\)
such that \(e\) codes a formula \(\varphi\) and \((\mathcal N,<)\) satisfies
\(\varphi\) under the assignment coded by \(s\) with \(c\) read as \(n\).
A sentence does not depend on the assignment, so \(g(e,n)\in U^\ast\) iff
\((\mathcal N,<,n)\models\psi_e\).

\emph{The sentence \(\sigma\).}  Let \(\chi\) be a sentence of
\(\mathscr L[\tau_<]\) whose models are exactly the copies of
\((\mathcal N,<)\): \(\chi_{\mathcal N}\) of
\cref{thm:p03-w03-p03-047-q0-not-re} for \(\mathscr L(Q_0)\), the sentence of
part~3 of \cref{thm:p03-w03-p03-047-weak-sol} for \(\mathscr L^{w2}\).  Let
\(\mathrm{Fml}\) define the set of formula codes.  Let \(\sigma\) be the
conjunction of \(\chi\) with the universal closures of
\begin{align*}
 &Ux\to\exists e\,\exists s\,\exists n\,\exists w\,(w=\pi(s,n)\land x=\pi(e,w)\land\mathrm{Fml}(e)),\\
 &\mathrm{At}(e)\to\bigl(U\langle e,s,n\rangle\leftrightarrow\mathrm{True}_{\mathrm{at}}(e,s,n)\bigr),\\
 &\mathrm{Neg}(e,e')\to\bigl(U\langle e,s,n\rangle\leftrightarrow\neg U\langle e',s,n\rangle\bigr),\\
 &\mathrm{Conj}(e,e_1,e_2)\to\bigl(U\langle e,s,n\rangle\leftrightarrow(U\langle e_1,s,n\rangle\land U\langle e_2,s,n\rangle)\bigr),\\
 &\mathrm{Ex}(e,i,e')\to\bigl(U\langle e,s,n\rangle\leftrightarrow\exists m\,\exists s'\,(\mathrm{Upd}(s,i,m,s')\land U\langle e',s',n\rangle)\bigr),
\end{align*}
and, for \(\mathscr L(Q_0)\),
\(\mathrm{Qz}(e,i,e')\to(U\langle e,s,n\rangle\leftrightarrow Q_0m\,\exists s'\,(\mathrm{Upd}(s,i,m,s')\land U\langle e',s',n\rangle))\),
or, for \(\mathscr L^{w2}\), where \(\mathrm{At}\) and
\(\mathrm{True}_{\mathrm{at}}\) also cover the atomic formulas
\(X^k_jt_1\cdots t_k\),
\(\mathrm{Ex2}(e,k,j,e')\to(U\langle e,s,n\rangle\leftrightarrow\exists r\,\exists s'\,(\mathrm{Upd2}(s,k,j,r,s')\land U\langle e',s',n\rangle))\).
Each conjunct is a sentence of \(\mathscr L[\tau_<\cup\{U\}]\): \(\chi\) is,
the clauses are first-order except for the single \(Q_0\), and both logics
contain first-order logic and are closed under \(\land\).

\emph{\((\mathcal N,<,U^\ast)\models\sigma\).}  \(\chi\) holds in
\((\mathcal N,<)\).  The first clause holds by the definition of
\(U^\ast\).  The clauses for atomic formulas, \(\neg\) and \(\land\) are
the satisfaction clauses for these cases.  For \(\exists v_i\): the formula
coded by \(e\) is satisfied under \(s\) iff some \(m\) makes the formula
coded by \(e'\) true under \(s\) updated at \(v_i\) to \(m\), and
\(\mathrm{Upd}\) defines that update.  For \(Q_0v_i\): the formula is
satisfied iff infinitely many \(m\) make the formula coded by \(e'\) true
under the updated assignment; since \(\mathrm{Upd}\) is the graph of a
function of \((s,i,m)\), the set of \(m\) with
\(\exists s'(\mathrm{Upd}(s,i,m,s')\land U^\ast\langle e',s',n\rangle)\) is
exactly that set, and \(Q_0m\) says it is infinite.  For \(\exists X^k_j\):
the formula is satisfied iff some finite \(k\)-ary relation on \(\omega\)
makes the body true; every such relation has a code \(r\) and every \(r\)
codes such a relation, so the clause holds.

\emph{\(U\) is implicitly defined by \(\sigma\).}  Let \((A,U_1)\) and
\((A,U_2)\) be expansions of one \(\tau_<\)-structure \(A\) that satisfy
\(\sigma\).  Then \(A\models\chi\), so there is an isomorphism
\(h:A\to(\mathcal N,<)\), and by the isomorphism property
\((\mathcal N,<,h[U_l])\models\sigma\) for \(l=1,2\).  It suffices to show
that every \(V\subseteq\omega\) with \((\mathcal N,<,V)\models\sigma\)
equals \(U^\ast\); then \(h[U_1]=h[U_2]\), and \(U_1=U_2\) since \(h\) is
injective.  By the first clause both \(V\) and \(U^\ast\) consist of
numbers \(\pi(e,\pi(s,n))\) with \(e\) a formula code.  We show by
induction on the formula \(\varphi\) coded by \(e\) that for all \(s,n\),
\(\pi(e,\pi(s,n))\in V\) iff \(\pi(e,\pi(s,n))\in U^\ast\).  For atomic
\(\varphi\) both sides are equivalent to \(\mathrm{True}_{\mathrm{at}}(e,s,n)\).
For \(\neg\) and \(\land\) the clauses reduce membership to membership for
the immediate subformulas, where the induction hypothesis applies.  For
\(\exists v_i\varphi'\), \(Q_0v_i\varphi'\) and \(\exists X^k_j\varphi'\),
the set of \(m\) (or \(r\)) satisfying the right-hand side of the clause is
the same for \(V\) and for \(U^\ast\) by the induction hypothesis for
\(\varphi'\), so the left-hand sides agree.  Hence \(V=U^\ast\).

Both logics are one-sorted logics with the atom, Boolean and
particularization properties: \(\neg\), \(\land\) and \(\exists\) are among
their formation rules, and a constant \(d\) is particularized by replacing
it with a variable not occurring in the sentence and quantifying that
variable existentially.  So part~(1) applies to both.
\end{proof}
\end{annexitem}

\begin{annexitem}{annex:p03-w03-p03-047-sol-failures}{thm:p03-w03-p03-047-sol-failures}{Infinity and uncountability in second-order logic}
\begin{proof}
Let \(\operatorname{Fn}(G)\) be the first-order formula in a binary relation
variable \(G\) saying that \(G\) is the graph of a function from the domain
to itself, \(\forall x\exists y\,(Gxy\land\forall y'(Gxy'\to y'=y))\), and let
\[
 \mathrm{Inf}:=\exists G\bigl(\operatorname{Fn}(G)\land
 \forall x\forall x'\forall y(Gxy\land Gx'y\to x=x')\land
 \exists z\forall x\,\neg Gxz\bigr).
\]
If the domain \(A\) is infinite, choose (by the Axiom of Choice) an injection
\(g:\omega\to A\) and let \(f(g(n))=g(n+1)\) and \(f(a)=a\) for \(a\) outside
the range of \(g\); \(f\) is injective and \(g(0)\) is not in its range, so
its graph witnesses \(\mathrm{Inf}\).  If \(A\) is finite, an injective map
\(A\to A\) is surjective by counting, so \(\mathrm{Inf}\) fails.

For uncountability, let \(\mathrm{Lin}(R)\) say that the binary relation
variable \(R\) is a strict linear order of the domain, and let
\(\mathrm{Inf}_R(y)\) be \(\mathrm{Inf}\) with each first-order quantifier
relativized to \(\{x:Rxy\}\) (and the clause \(\operatorname{Fn}\) asserting
a map of that set to itself).  By the previous paragraph applied to the set
of \(R\)-predecessors of \(y\) (when that set is empty, the relativized
\(\exists z\) clause fails, matching finiteness), \(\mathrm{Inf}_R(y)\) holds iff \(y\) has
infinitely many \(R\)-predecessors.  Put
\[
 \mathrm{Unc}:=\neg\exists R\bigl(\mathrm{Lin}(R)\land\forall y\,\neg
 \mathrm{Inf}_R(y)\bigr).
\]
If \(A\) has power at most \(\aleph_0\), fix an injection
\(e:A\to\omega\) and let \(a\mathrel{R}b\) iff \(e(a)<e(b)\); every \(b\) has
at most \(e(b)\) predecessors, so the existential holds and \(\mathrm{Unc}\)
fails.  If such an \(R\) exists, the map \(a\mapsto\) number of
\(R\)-predecessors of \(a\) is injective, by the argument of
\cref{annex:p03-w03-p03-047-q0-omega}, so \(A\) has power at most
\(\aleph_0\).  Hence \(\mathrm{Unc}\) holds exactly in the uncountable
structures.

For \(n\ge1\) let \(\sigma'_n\) be the first-order sentence
\(\exists x_1\ldots x_n\bigwedge_{i<j\le n}x_i\ne x_j\).  The set
\(\{\neg\mathrm{Inf}\}\cup\{\sigma'_n:n\ge1\}\) is countable, finitely
satisfiable in finite structures and has no model, so countable compactness
fails.  The sentence \(\mathrm{Unc}\) has the infinite model \(\omega_1\) and
no model of power at most \(\aleph_0\), so the L\"owenheim property fails.
\end{proof}
\end{annexitem}

\begin{annexitem}{annex:p03-w03-p03-047-lindstrom-bridge}{thm:p03-w03-p03-047-lindstrom-bridge}{Matching the two forms of Lindstr\"om's hypotheses}
\begin{proof}
(1) Barwise's form says: if a countable set has no model, some finite subset
has no model.  Hypothesis~1 of \cref{thm:p03-lindstrom} says: if every
finite subset of a countable set has a model, the set has a model.  Each is
the contrapositive of the other.

(2) Assume Barwise's property and let \(\varphi\) be satisfiable, with model
\(\mathfrak M\).  If \(\mathfrak M\) is finite it has power at most
\(\aleph_0\).  If \(\mathfrak M\) is infinite, the property gives a model of
power at most \(\aleph_0\).  So hypothesis~2 holds.  Conversely, if every
satisfiable sentence has a model of power at most \(\aleph_0\), then in
particular so does every sentence with an infinite model, which is Barwise's
property under the reading of ``countable'' fixed in the body.  The final
assertion is \cref{thm:p03-lindstrom} applied to these hypotheses.
\end{proof}
\end{annexitem}

\begin{annexitem}{annex:p03-w03-p03-047-craig-separation}{thm:p03-w03-p03-047-craig-separation}{Interpolation as separation of PC classes}
\begin{proof}
(1), interpolation implies separation.  Let \(K_0,K_1\subseteq
\operatorname{Str}[\tau]\) be disjoint, \(K_i\) being PC via
\(\varphi_i\in\mathscr L[\tau_i']\) with \(\tau\subseteq\tau_i'\).  First
make the added symbols disjoint.  Choose a renaming \(\rho\) of \(\tau_1'\)
that is the identity on \(\tau\) and sends each symbol of
\(\tau_1'\setminus\tau\) to a fresh symbol of the same kind and arity outside
\(\tau_0'\cup\tau_1'\).  By the renaming property there is \(\varphi_1^\rho\)
with \(\mathfrak N\models\varphi_1\) iff \(\mathfrak N^\rho\models\varphi_1^\rho\).
Because \(\rho\) fixes \(\tau\), \(\mathfrak N\mapsto\mathfrak N^\rho\) is a
bijection between the expansions of a \(\tau\)-structure \(\mathfrak M\) to
\(\tau_1'\) and its expansions to \(\rho[\tau_1']\); so \(K_1\) is PC via
\(\varphi_1^\rho\).  Replacing \(\varphi_1\) by \(\varphi_1^\rho\) we may
assume \(\tau_0'\cap\tau_1'=\tau\).

Then \(\varphi_0\models\neg\varphi_1\).  Indeed, let \(\mathfrak B\) be a
\((\tau_0'\cup\tau_1')\)-structure with \(\mathfrak B\models\varphi_0\).  By
the reduct property \(\mathfrak B\restriction\tau_0'\models\varphi_0\), so
\(\mathfrak B\restriction\tau\in K_0\), hence
\(\mathfrak B\restriction\tau\notin K_1\).  As \(\mathfrak B\restriction\tau_1'\)
is an expansion of \(\mathfrak B\restriction\tau\), it does not satisfy
\(\varphi_1\); so \(\mathfrak B\not\models\varphi_1\), that is,
\(\mathfrak B\models\neg\varphi_1\).  Interpolation gives
\(\theta\in\mathscr L[\tau_0'\cap\tau_1']=\mathscr L[\tau]\) with
\(\varphi_0\models\theta\) and \(\theta\models\neg\varphi_1\).  If
\(\mathfrak M\in K_0\), take an expansion \(\mathfrak N\models\varphi_0\) to
\(\tau_0'\) and expand it arbitrarily to \(\tau_0'\cup\tau_1'\); this
structure satisfies \(\varphi_0\), hence \(\theta\), and by the reduct
property \(\mathfrak M\models\theta\).  If \(\mathfrak M\in K_1\) and
\(\mathfrak M\models\theta\), take an expansion \(\mathfrak N\models\varphi_1\)
to \(\tau_1'\) and expand it to \(\mathfrak B\) over
\(\tau_0'\cup\tau_1'\); then \(\mathfrak B\models\theta\), so
\(\mathfrak B\models\neg\varphi_1\), while
\(\mathfrak B\restriction\tau_1'=\mathfrak N\models\varphi_1\), a
contradiction.  So \(\theta\) separates \(K_0\) from \(K_1\).

(1), separation implies interpolation.  Let \(\varphi_0\in\mathscr L[\tau_0]\),
\(\varphi_1\in\mathscr L[\tau_1]\), \(\varphi_0\models\varphi_1\), and
\(\tau=\tau_0\cap\tau_1\).  Let \(K_0\) be the class of reducts to \(\tau\) of
models of \(\varphi_0\), and \(K_1\) that of reducts of models of
\(\neg\varphi_1\in\mathscr L[\tau_1]\); both are PC.  They are disjoint: if
\(\mathfrak M=\mathfrak B_0\restriction\tau=\mathfrak B_1\restriction\tau\)
with \(\mathfrak B_0\models\varphi_0\) over \(\tau_0\) and
\(\mathfrak B_1\models\neg\varphi_1\) over \(\tau_1\), then since
\(\tau_0\setminus\tau\) and \(\tau_1\setminus\tau\) are disjoint there is a
\((\tau_0\cup\tau_1)\)-structure \(\mathfrak C\) with
\(\mathfrak C\restriction\tau_i=\mathfrak B_i\).  By the reduct property
\(\mathfrak C\models\varphi_0\), so \(\mathfrak C\models\varphi_1\), so
\(\mathfrak B_1\models\varphi_1\), contradicting
\(\mathfrak B_1\models\neg\varphi_1\).  Let \(\theta\in\mathscr L[\tau]\)
separate them.  If \(\mathfrak B\models\varphi_0\) then
\(\mathfrak B\restriction\tau\in K_0\), so \(\mathfrak B\models\theta\).  If
\(\mathfrak B\models\theta\) but \(\mathfrak B\not\models\varphi_1\), then
\(\mathfrak B\restriction\tau_1\models\neg\varphi_1\), so
\(\mathfrak B\restriction\tau\in K_1\) satisfies \(\theta\), which separation
excludes.  Hence \(\varphi_0\models\theta\models\varphi_1\).

(2) Let \(K\) be PC and co-PC in \(\tau\).  Then \(K\) and its complement are
disjoint PC classes; by~(1) some \(\theta\in\mathscr L[\tau]\) has
\(K\subseteq\operatorname{Mod}(\theta)\) and
\((\operatorname{Str}[\tau]\setminus K)\cap\operatorname{Mod}(\theta)=\varnothing\),
that is, \(\operatorname{Mod}(\theta)=K\).
\end{proof}
\end{annexitem}

\begin{annexitem}{annex:p03-w03-p03-047-bounded-omega}{thm:p03-w03-p03-047-bounded-omega}{Countable compactness bounds well-ordering}
\begin{proof}
\emph{First-order sentences.}  Every first-order sentence \(\chi\) of a
vocabulary \(\tau\) has an \(\mathscr L\)-sentence of \(\tau\) with the same
models.  Induct on \(\chi\), simultaneously for all vocabularies.  Atomic
sentences are covered by the atom property, \(\neg\) and \(\land\) by the
Boolean property.  For \(\exists x\,\chi_0(x)\), take a constant \(c\) not in
\(\tau\); the sentence \(\chi_0(c)\) of \(\tau\cup\{c\}\) is shorter, so it has
an equivalent \(\psi\in\mathscr L[\tau\cup\{c\}]\), and the
particularization property gives \(\psi'\in\mathscr L[\tau]\) with
\(\mathfrak A\models\psi'\) iff \((\mathfrak A,a)\models\psi\) for some
\(a\), iff \((\mathfrak A,a)\models\chi_0(c)\) for some \(a\), iff
\(\mathfrak A\models\exists x\,\chi_0\).  The other connectives and
\(\forall\) are expressed by these.

\emph{Bound.}  Let \(\varphi\in\mathscr L[\tau]\) pin down ordinals via
\(<\), and suppose that for every \(m\) some model of \(\varphi\) has
\(<\) a well-ordering of its field of order type at least \(m+1\) (this
holds if \(\varphi\) pins down an infinite ordinal or arbitrarily large
finite ones).  Add new constants \(c_n\) (\(n\in\omega\)) and let
\(T=\{\varphi\}\cup\{c_{n+1}<c_n:n\in\omega\}\), a countable set of
\(\mathscr L\)-sentences by property~(i) of
\cref{def:p03-w03-p03-047-ebbinghaus-logic} and the atom property.  A finite
subset of \(T\) is contained in \(\{\varphi\}\cup\{c_{n+1}<c_n:n<m\}\) for
some \(m\).  Take a model \(\mathfrak A\) of \(\varphi\) whose field has
order type at least \(m+1\), with first elements \(e_0<\cdots<e_m\), and
interpret \(c_k\) as \(e_{m-k}\) for \(k\le m\) and as \(e_0\) otherwise.
The constraints \(c_{n+1}<c_n\) for \(n<m\) hold, and the expansion still
satisfies \(\varphi\) by the reduct property.  By countable compactness
\(T\) has a model \(\mathfrak B\).  Its reduct to \(\tau\) satisfies
\(\varphi\), so \(<^{\mathfrak B}\) well-orders its field; but the
\(c_n^{\mathfrak B}\) lie in the field and form a strictly descending
sequence.  This contradiction shows that the ordinals pinned down by
\(\varphi\) are bounded by some natural number.

\emph{Finite ordinals.}  For \(n\ge2\) the first-order sentence saying that
\(<\) is a strict linear order whose field
\(\{x:\exists y(x<y\lor y<x)\}\) has exactly \(n\) elements has a model of
order type \(n\) (the order on \(\{0,\ldots,n-1\}\)), and in every model the
field is a finite linear order, hence a well-ordering; by the first
paragraph it has an \(\mathscr L\)-equivalent, which pins down \(n\).

So no sentence pins down arbitrarily large ordinals, \(\mathscr L\) is
bounded, every ordinal pinned down by a sentence is finite, and arbitrarily
large finite ordinals are pinned down; the supremum \(w(\mathscr L)\) is
\(\omega\).  The case \(\mathscr L_{\omega\omega}\) uses
\cref{thm:p02-w02-p02-005-bw-compactness}; the case \(\mathscr L(Q_1)\) uses its
countable compactness from \cref{thm:p03-keisler-q-one} and the fact that it
contains first-order logic.
\end{proof}
\end{annexitem}

\begin{annexitem}{annex:p03-w03-p03-047-hanf-exists}{thm:p03-w03-p03-047-hanf-exists}{Existence of the Hanf number}
\begin{proof}
Let \(S\) be the set of sentences of \(\mathscr L\) and \(S'\) the set of
those that have a model but do not have models of arbitrarily large power.
For \(\varphi\in S'\) there is a cardinal above the power of every model of
\(\varphi\); let \(\mu_\varphi\) be the least such cardinal.  By Replacement
\(\{\mu_\varphi:\varphi\in S'\}\) is a set; let \(\kappa_0\) be its supremum
(\(\kappa_0=0\) if \(S'=\varnothing\)).  If \(\varphi\in S\) has a model of
power at least \(\kappa_0\), then \(\varphi\notin S'\), since every model of
a member of \(S'\) has power below \(\mu_\varphi\le\kappa_0\); having a model,
\(\varphi\) therefore has models of arbitrarily large power.  So \(\kappa_0\)
has the defining property, and the least cardinal with that property exists.

For first-order logic in a countable vocabulary the sentences form a
countable set.  A sentence with a model of power at least \(\aleph_0\) has
an infinite model \(\mathfrak M\); for every cardinal
\(\kappa\ge\max(\lvert\dom\mathfrak M\rvert,\aleph_0)\),
\cref{thm:p02-w02-p02-033-upward-ls-inclusion} gives an elementary
extension of \(\mathfrak M\) of power \(\kappa\), which satisfies the
sentence.  So the sentence has arbitrarily large models, and \(\aleph_0\)
has the property.  No finite \(n\ge1\)
has it: the sentence \(\forall x_0\ldots x_n\bigvee_{i<j\le n}x_i=x_j\) has a
model of power \(n\) and no model of power above \(n\); its instance for
\(n=1\) also excludes the cardinal \(0\).  Hence
\(h(\mathscr L_{\omega\omega})=\aleph_0\).
\end{proof}
\end{annexitem}

\begin{annexitem}{annex:p03-w03-p03-047-omega-logic-renaming}{thm:p03-w03-p03-047-omega-logic-renaming}{\(\omega\)-logic and the renaming property}
\begin{proof}
(1) Property (i): if \(\tau\subseteq\sigma\) and \(\{s,<\}\subseteq\tau\),
then \(\{s,<\}\subseteq\sigma\) and first-order sentences of \(\tau\) are
first-order sentences of \(\sigma\); otherwise \(\mathscr L[\tau]\) is empty.
Property (ii) is built into the satisfaction clause.  Property (iii): an
isomorphism preserves first-order satisfaction and maps
\(\mathfrak A\restriction\{s,<\}\) isomorphically onto
\(\mathfrak B\restriction\{s,<\}\).  Property (iv): if
\(\varphi\in\mathscr L[\tau]\) and \(\tau\subseteq\tau_{\mathfrak A}\), then
\(\{s,<\}\subseteq\tau\), \(\varphi\) lies in both
\(\mathscr L[\tau]\) and \(\mathscr L[\tau_{\mathfrak A}]\),
\((\mathfrak A\restriction\tau)\restriction\{s,<\}=\mathfrak A\restriction\{s,<\}\),
and first-order satisfaction of \(\varphi\) depends only on the reduct to
\(\tau\).  Property (v) fails for \(\rho:\{s,<\}\to\{s,\prec\}\): the sentence
\(\forall x\,x=x\) lies in \(\mathscr L[\{s,<\}]\), but
\(\mathscr L[\{s,\prec\}]=\varnothing\) because \(<\notin\{s,\prec\}\), so no
\(\varphi^\rho\) exists.

(2) Suppose \(\psi\) is a first-order sentence of \(\{s,\prec\}\) with
\(\mathfrak A\models\forall x\,x=x\) iff \(\mathfrak A^\rho\models\psi\) for
all \(\{s,<\}\)-structures \(\mathfrak A\).  The left side holds iff
\(\mathfrak A\cong(\omega,<)\).  Every \(\{s,\prec\}\)-structure is
\(\mathfrak A^\rho\) for exactly one \(\mathfrak A\), and
\(\mathfrak A\cong(\omega,<)\) iff \(\mathfrak A^\rho\cong(\omega,\prec)\)
with \(\prec\) read as the usual order.  So \(\psi\) holds exactly in the
copies of \((\omega,<)\) written with \(\prec\).  That structure is infinite,
and by \cref{thm:p02-w02-p02-033-upward-ls-inclusion} with
\(\kappa=\aleph_1\) it has an elementary extension of power \(\aleph_1\).
That extension satisfies \(\psi\) and is not a copy of \((\omega,<)\).  This contradiction proves~(2).
\end{proof}
\end{annexitem}

\begin{annexitem}{annex:p03-w03-p03-047-function-elimination}{thm:p03-w03-p03-047-function-elimination}{Graph\\elimination}
\begin{proof}
Regard \(\varphi\) as a sentence over \(\tau\cup\sigma\), using
vocabulary extension and the reduct property. Simultaneously replace
each function \(f\) by the graph predicate \(R_f(\bar x,y)\),
and each constant \(c\) by \(R_c(y)\), leaving the shared relational
symbols unchanged. The substitution property supplies a sentence
\(\varphi^{\mathrm{gr}}\) over \(\sigma\).

A \(\sigma\)-structure \(\mathfrak B\) satisfies \(\Gamma_\tau\)
if and only if each proposed graph is total and single-valued, and
each proposed constant predicate has exactly one element. In that
case define \(f^{\mathfrak A}(\bar b)\) to be the unique \(b\)
with \(R_f^{\mathfrak B}(\bar b,b)\), and define
\(c^{\mathfrak A}\) to be the unique member of
\(R_c^{\mathfrak B}\). Keep the same domain and the same original
relations. This determines a unique \(\tau\)-structure
\(\mathfrak A\), and \(\mathfrak B=\mathfrak A^{\mathrm{gr}}\).
All substitution interpretations are now admissible. Hence the
substitution equivalence gives
\[
 \mathfrak B\models\varphi^{\mathrm{gr}}
 \quad\Longleftrightarrow\quad
 \mathfrak A\models\varphi.
\]
This proves the first assertion and both directions of the second.
It imposes no condition on the translated sentence's value at
inadmissible structures, because \(\Gamma_\tau\) already excludes them.
Each member of \(\Gamma_\tau\) is first-order and therefore available
in the regular logic. If the set of eliminated function and constant
symbols is finite, so is \(\Gamma_\tau\); its finite conjunction
with \(\varphi^{\mathrm{gr}}\) defines exactly the required class
over all structures. For the empty family use a first-order tautology
as the empty conjunction. This proves the last assertion.
\end{proof}
\end{annexitem}

\begin{annexitem}{annex:p03.graph.1}{thm:p03-finite-occurrence-graph-obstruction}{Graph admissibility}
\label{annex:p03.finite-occurrence-graph-obstruction}
\begin{proof}
Suppose a sentence \(\theta\) defined the stated class. By finite
occurrence there is a finite sub-vocabulary \(\sigma_0\) such that
\(\theta\in\mathscr L[\sigma_0]\). Choose \(m\) with
\(R_m\notin\sigma_0\). On a one-element domain \(\{a\}\), let
\(\mathfrak A\) interpret every \(R_n\) as \(\{(a,a)\}\).
Every relation is then the graph of a total unary function, so
\(\mathfrak A\models\theta\). Form \(\mathfrak B\) on the same
domain, changing only \(R_m\) to the empty relation.
The two structures have identical \(\sigma_0\)-reducts, so the reduct
property implies \(\mathfrak B\models\theta\). But \(R_m^{\mathfrak B}\)
is not total, contrary to the defining property of \(\theta\).
Thus no such sentence exists. First-order sentences have finite
occurrence and the reduct property, giving the stated special case.
\end{proof}
\end{annexitem}

\begin{annexitem}{annex:p03-w03-p03-047-substitution-renaming}{thm:p03-w03-p03-047-substitution-renaming}{Renaming obtained by substitution}
\begin{proof}
\emph{Disjoint case.}  Let \(\eta:\alpha\to\beta\) be a renaming with
\(\alpha\cap\beta=\varnothing\), and \(\varphi\in\mathscr L[\alpha]\).  By
property~(i), \(\varphi\in\mathscr L[\alpha\cup\beta]\); apply substitution
with smaller vocabulary \(\beta\), so that every symbol of \(\alpha\) is
replaced.  Use the atom property to choose: for an \(n\)-ary relation
\(R\), \(\psi_R(\mathbf c)\) equivalent to \(\eta(R)\mathbf c\); for an
\(n\)-ary function \(f\), \(\psi_f(\mathbf c,c)\) equivalent to
\(\eta(f)(\mathbf c)=c\); for a constant \(d\), \(\psi_d(e)\) equivalent to
\(\eta(d)=e\).  Let \(\varphi^\eta=\varphi^\ast\in\mathscr L[\beta]\).  For a
\(\beta\)-structure \(\mathfrak B\) the defined sets are always graphs, and
the required expansion \(\mathfrak B'\) to \(\alpha\cup\beta\) is unique:
\(S^{\mathfrak B'}=\eta(S)^{\mathfrak B}\) for \(S\in\alpha\).  Its reduct to
\(\alpha\) is \(\mathfrak B^{\eta^{-1}}\).  So \(\mathfrak B\models\varphi^\eta\)
iff \(\mathfrak B'\models\varphi\) iff \(\mathfrak B^{\eta^{-1}}\models\varphi\)
by the reduct property.  For \(\mathfrak A\in\operatorname{Str}[\alpha]\)
and \(\mathfrak B=\mathfrak A^\eta\) this reads
\(\mathfrak A\models\varphi\) iff \(\mathfrak A^\eta\models\varphi^\eta\).

\emph{General case.}  Let \(\rho:\tau\to\sigma\) be a renaming.  Using the
unbounded supply of symbols, choose a vocabulary \(\pi\) disjoint from
\(\tau\cup\sigma\) and a renaming \(\rho_1:\tau\to\pi\); then
\(\rho_2=\rho\circ\rho_1^{-1}:\pi\to\sigma\) is a renaming between disjoint
vocabularies, and \(\mathfrak A^\rho=(\mathfrak A^{\rho_1})^{\rho_2}\).  The
disjoint case gives \(\varphi^{\rho_1}\in\mathscr L[\pi]\) and then
\((\varphi^{\rho_1})^{\rho_2}\in\mathscr L[\sigma]\), with
\(\mathfrak A\models\varphi\) iff \(\mathfrak A^{\rho_1}\models\varphi^{\rho_1}\)
iff \(\mathfrak A^\rho\models(\varphi^{\rho_1})^{\rho_2}\).  Only one sort
occurs, so no sort symbol has to be renamed.
\end{proof}
\end{annexitem}

\begin{annexitem}{annex:p03-w03-p03-047-craig-beth}{thm:p03-w03-p03-047-craig-beth}{Beth's property from interpolation}
\begin{proof}
Let \(R\in\tau\) be \(n\)-ary and implicitly defined by \(\varphi\).  Choose
a fresh \(n\)-ary \(R'\), put \(\tau'=(\tau\setminus\{R\})\cup\{R'\}\), and let
\(\rho:\tau\to\tau'\) be the renaming sending \(R\) to \(R'\) and fixing the
rest; let \(\varphi'=\varphi^\rho\).  Let \(\mathbf c=c_1\ldots c_n\) be
fresh constants.  Using properties (i), the atom property and the Boolean
property, let \(\chi_0\in\mathscr L[\tau\cup\{\mathbf c\}]\) be
\(\varphi\land R\mathbf c\) and \(\chi_1\in\mathscr L[\tau'\cup\{\mathbf c\}]\)
be \(\neg(\varphi'\land\neg R'\mathbf c)\).

\emph{Claim: \(\chi_0\models\chi_1\).}  Let \(\mathfrak C\) be a
\((\tau\cup\tau'\cup\{\mathbf c\})\)-structure with
\(\mathfrak C\models\varphi\land R\mathbf c\) and suppose
\(\mathfrak C\models\varphi'\).  Let \(\mathfrak D=\mathfrak C\restriction
(\tau\setminus\{R\})\).  Two expansions of \(\mathfrak D\) to \(\tau\)
satisfy \(\varphi\): \(\mathfrak E_1=\mathfrak C\restriction\tau\) by the
reduct property, and \(\mathfrak E_2=(\mathfrak C\restriction\tau')^{\rho^{-1}}\),
which interprets \(R\) as \(R'^{\mathfrak C}\), because
\(\mathfrak E_2^\rho=\mathfrak C\restriction\tau'\models\varphi^\rho\).  By
implicit definability \(\mathfrak E_1=\mathfrak E_2\), so
\(R^{\mathfrak C}=R'^{\mathfrak C}\).  Since
\(\mathbf c^{\mathfrak C}\in R^{\mathfrak C}\), \(\mathfrak C\models R'\mathbf c\),
and \(\mathfrak C\models\chi_1\).

Interpolation yields \(\psi\) in
\(\mathscr L[(\tau\cup\{\mathbf c\})\cap(\tau'\cup\{\mathbf c\})]=
\mathscr L[(\tau\setminus\{R\})\cup\{\mathbf c\}]\) with
\(\chi_0\models\psi\models\chi_1\).  Let \(\mathfrak A\models\varphi\) be a
\(\tau\)-structure and \(\mathbf a\in A^n\).  If
\(\mathbf a\in R^{\mathfrak A}\), then \((\mathfrak A,\mathbf a)\models\chi_0\),
so \((\mathfrak A,\mathbf a)\models\psi\).  If \((\mathfrak A,\mathbf a)\models\psi\),
let \(\mathfrak C\) expand \((\mathfrak A,\mathbf a)\) by
\(R'^{\mathfrak C}=R^{\mathfrak A}\).  Then
\(\mathfrak C\restriction\tau'=\mathfrak A^\rho\models\varphi'\) and
\(\mathfrak C\models\psi\), hence \(\mathfrak C\models\chi_1\); with
\(\mathfrak C\models\varphi'\) this gives \(\mathfrak C\models R'\mathbf c\),
that is, \(\mathbf a\in R^{\mathfrak A}\).  Therefore
\(R^{\mathfrak A}=\{\mathbf a:(\mathfrak A,\mathbf a)\models\psi(\mathbf c)\}\)
for every model \(\mathfrak A\) of \(\varphi\).  A regular logic has the
properties used, so it has the Beth property for relation symbols.
\end{proof}
\end{annexitem}

\begin{annexitem}{annex:p03-w03-p03-047-q1-not-aleph1}{thm:p03-w03-p03-047-q1-not-aleph1}{The \(\aleph_1\)-compactness counterexample}
\begin{proof}
The set has one sentence for each pair \(\alpha<\beta<\aleph_1\) plus one,
so its power is \(\aleph_1\).  A finite subset mentions finitely many
constants \(c_{\alpha_1},\ldots,c_{\alpha_k}\).  Interpret them in the domain
\(\omega\) as \(1,\ldots,k\) and every other constant as \(0\).  The
mentioned constants are pairwise distinct, so every inequality of the subset
holds, and the countable domain satisfies \(\neg Q_1x\,(x=x)\).  A model of
the whole set would contain the \(\aleph_1\) pairwise distinct elements
\(c_\alpha\), so \(Q_1x\,(x=x)\) would hold in it, contradicting the first
sentence.
\end{proof}
\end{annexitem}

\begin{annexitem}{annex:p03-w03-p03-047-sol-interpolation}{thm:p03-w03-p03-047-sol-interpolation}{Quantifying away the private symbols}
\begin{proof}
The sentence \(\varphi_0\) contains finitely many nonlogical symbols; let
\(S_1,\ldots,S_k\) be those not in \(\tau_1\).  Choose distinct second-order
variables \(X_1,\ldots,X_k\) not occurring in \(\varphi_0\), each of the kind
and arity of the corresponding \(S_i\) (a relation, a function, or for a
constant an individual variable), and let
\(\psi=\exists X_1\ldots\exists X_k\,\varphi_0[S_1/X_1,\ldots,S_k/X_k]\).
Every nonlogical symbol of \(\psi\) occurs in \(\varphi_0\) and is in
\(\tau_1\), so \(\psi\in\mathscr L^2[\tau_0\cap\tau_1]\).

\emph{Substitution lemma.}  For a structure \(\mathfrak B\) interpreting the
symbols of \(\varphi_0\) other than \(S_1,\ldots,S_k\), objects
\(V_1,\ldots,V_k\) of the right kinds, and any assignment \(s\):
\(\mathfrak B,s[X_i\mapsto V_i]\models\varphi_0[S/X]\) iff
\(\mathfrak B[S_i\mapsto V_i],s\models\varphi_0\), where
\(\mathfrak B[S_i\mapsto V_i]\) interprets each \(S_i\) as \(V_i\).  By
induction on subformulas: atomic formulas and terms receive the same values
because \(X_i\) and \(S_i\) are interpreted by the same object; the Boolean
cases are immediate; and quantifier cases (first or second order) commute
with the change because the \(X_i\) are not bound inside \(\varphi_0\).

If \(\mathfrak B\models\varphi_0\), take \(V_i=S_i^{\mathfrak B}\); by the
lemma \(\mathfrak B\models\psi\).  If \(\mathfrak B\models\psi\) for a
\((\tau_0\cup\tau_1)\)-structure, choose witnesses \(V_i\); by the lemma
\(\mathfrak B[S_i\mapsto V_i]\models\varphi_0\), hence it satisfies
\(\varphi_1\).  The \(S_i\) are not in \(\tau_1\), so
\(\mathfrak B[S_i\mapsto V_i]\restriction\tau_1=\mathfrak B\restriction\tau_1\),
and \(\mathfrak B\models\varphi_1\).  Thus
\(\varphi_0\models\psi\models\varphi_1\).  The single domain is shared by
\(\mathfrak B\) and \(\mathfrak B[S_i\mapsto V_i]\); this is the step that has
no counterpart for a sort occurring only in \(\tau_0\).
\end{proof}
\end{annexitem}

\begin{annexitem}{annex:p03-w03-p03-047-sol-ch}{thm:p03-w03-p03-047-sol-ch}{A second-order sentence for the continuum hypothesis}
\begin{proof}
All relation variables below range over arbitrary relations on the domain
\(D\), as in \(\mathscr L^2\); \(\varphi\to\psi\) abbreviates
\(\neg\varphi\lor\psi\).

\emph{Copies of \(\omega\).}  Let \(\Omega(N,L)\), with \(N\) unary and \(L\)
binary, be the conjunction of: \(\forall x\forall y\,(Lxy\to Nx\land Ny)\);
\(L\) is irreflexive, transitive and total on \(N\); \(\exists x\,Nx\);
\(\forall x\,(Nx\to\exists y\,Lxy)\);
\(\forall Z\,[(\exists x\,Zx\land\forall x\,(Zx\to Nx))\to\exists x\,(Zx\land\forall y\,(Zy\to\neg Lyx))]\);
and \(\forall x\,(Nx\land\exists y\,Lyx\to\exists y\,(Lyx\land\neg\exists z\,(Lyz\land Lzx)))\).
Claim: \(\Omega(N,L)\) holds iff \((N,L)\cong(\omega,<)\).  The copies of
\((\omega,<)\) satisfy each clause.  Conversely, by the fifth clause \(L\)
well-orders \(N\), so \((N,L)\cong(\alpha,<)\) for an ordinal \(\alpha\); by
the third clause \(\alpha\neq0\); by the fourth \(\alpha\) has no largest
element, so it is a limit ordinal; by the sixth every nonzero
\(\beta<\alpha\) is a successor.  If \(\alpha>\omega\), then \(\omega<\alpha\)
is nonzero and not a successor; so \(\alpha\le\omega\), and the only nonzero
limit ordinal \(\le\omega\) is \(\omega\).

\emph{Countable subsets.}  Let \(\mathrm{Ctbl}(Y)\) be
\(\exists N\exists L\exists F\,[\Omega(N,L)\land\forall x\,(Yx\to\exists y\,(Ny\land Fxy))\land\forall x\forall y\forall y'\,(Fxy\land Fxy'\to y=y')\land\forall x\forall x'\forall y\,(Yx\land Yx'\land Fxy\land Fx'y\to x=x')]\).
If it holds, choosing for each \(x\in Y\) the unique \(y\) with \(Fxy\)
gives an injection of \(Y\) into \(N\), which is countable, so \(Y\) is
countable.  If \(D\) is infinite and \(Y\subseteq D\) is countable, choose a
countably infinite \(N\subseteq D\) and a bijection \(g:N\to\omega\) (an
infinite set has a countably infinite subset, by the axiom of choice), put
\(L=\{(a,b):g(a)<g(b)\}\), and let \(F\) be the graph of an injection of
\(Y\) into \(N\); then \(\mathrm{Ctbl}(Y)\) holds.  So in structures with
infinite domain, \(\mathrm{Ctbl}(Y)\) holds iff \(Y\) is countable.

\emph{Equinumerosity with the domain.}  Let \(\mathrm{Eq}(Y)\) be
\(\exists F\,[\forall x\,(Yx\to\exists y\,Fxy)\land\forall x\forall y\forall y'\,(Fxy\land Fxy'\to y=y')\land\forall x\forall x'\forall y\,(Yx\land Yx'\land Fxy\land Fx'y\to x=x')\land\forall y\,\exists x\,(Yx\land Fxy)]\).
It holds iff \(F\restriction Y\) can be chosen as a bijection of \(Y\) onto
\(D\), that is, iff \(|Y|=|D|\).

\emph{Power of the continuum.}  Let \(\mathrm{Cont}\) be
\(\exists N\exists L\exists E\,[\Omega(N,L)\land\forall d\forall d'\,(\forall n\,(Nn\to(Edn\leftrightarrow Ed'n))\to d=d')\land\forall Z\,(\forall n\,(Zn\to Nn)\to\exists d\,\forall n\,(Nn\to(Edn\leftrightarrow Zn)))]\).
If it holds, \(d\mapsto\{n\in N:Edn\}\) is injective by the second clause
and onto \(\mathcal P(N)\) by the third, so \(|D|=|\mathcal P(N)|=2^{\aleph_0}\).
If \(|D|=2^{\aleph_0}\), \(D\) is infinite; choose \(N\), \(L\) as above and a
bijection \(g:D\to\mathcal P(N)\), and put \(E=\{(d,n):n\in g(d)\}\); then
\(\mathrm{Cont}\) holds.

\emph{The sentence.}  Let
\(\chi_{\mathrm{CH}}=\mathrm{Cont}\to\forall Y\,(\mathrm{Ctbl}(Y)\lor\mathrm{Eq}(Y))\).
Assume \(2^{\aleph_0}=\aleph_1\) and let \(D\) be any domain.  If
\(\mathrm{Cont}\) fails, \(\chi_{\mathrm{CH}}\) holds.  Otherwise
\(|D|=2^{\aleph_0}=\aleph_1\) and \(D\) is infinite.  For \(Y\subseteq D\):
if \(Y\) is countable, \(\mathrm{Ctbl}(Y)\) holds; if not,
\(\aleph_1\le|Y|\le|D|=\aleph_1\), so \(\mathrm{Eq}(Y)\) holds.  Thus
\(\chi_{\mathrm{CH}}\) is valid.  Assume \(2^{\aleph_0}\ne\aleph_1\), so
\(\aleph_1<2^{\aleph_0}\).  Take \(D=\mathcal P(\omega)\); \(\mathrm{Cont}\)
holds.  Choose \(Y\subseteq D\) with \(|Y|=\aleph_1\).  \(Y\) is uncountable
and \(D\) infinite, so \(\mathrm{Ctbl}(Y)\) fails; \(|Y|\ne|D|\), so
\(\mathrm{Eq}(Y)\) fails.  So \(\chi_{\mathrm{CH}}\) fails in \(D\).
\end{proof}
\end{annexitem}

\begin{annexitem}{annex:p03-w03-p03-047-weak-sol}{thm:p03-w03-p03-047-weak-sol}{Weak second-order logic}
\begin{proof}
(1) \(\exists X\,\forall x\,Xx\) holds iff some finite set contains the whole
domain, iff the domain is finite.  Likewise
\(\exists X\,\forall x\,(\varphi(x)\to Xx)\) holds (at given values of the
other free variables) iff the set \(\varphi^{\mathfrak A}\) of elements
satisfying \(\varphi\) is included in a finite set, iff
\(\varphi^{\mathfrak A}\) is finite; its negation holds iff
\(\varphi^{\mathfrak A}\) is infinite.  Define \(\varphi^w\) by recursion,
commuting with atomic formulas, connectives and \(\exists,\forall\), and
sending \(Q_0x\,\varphi\) to \(\neg\exists X\,\forall x\,(\varphi^w\to Xx)\)
with \(X\) fresh.  By induction on \(\varphi\), using the equivalence just
shown at the \(Q_0\) step, \(\varphi\) and \(\varphi^w\) are satisfied by the
same assignments in every structure.  The recursion is a syntactic
recursion on finite strings, hence computable under any coding in which the
syntactic operations are.

(2) Let \(\mathfrak G\) be a group and \(x\in G\).  If \(x\) has finite order
\(n\ge1\), \(X=\{e,x,\ldots,x^{n-1}\}\) is finite, contains \(e\), and is
closed under right multiplication by \(x\), since \(x^{n-1}x=e\).  If a
finite \(X\) contains \(e\) and is closed under right multiplication by
\(x\), then \(x^i\in X\) for all \(i\in\omega\) by induction on \(i\); as
\(X\) is finite, \(x^i=x^j\) for some \(i<j\), and \(x^{j-i}=e\).  So the
sentence holds iff every element has finite order.

(3) The sentence \((\chi_{\mathcal N})^w\), with \(\chi_{\mathcal N}\) from
\cref{thm:p03-w03-p03-047-q0-not-re}, has the same models as
\(\chi_{\mathcal N}\) by part~1, hence exactly the copies of
\((\mathcal N,<)\) by part~1 of \cref{thm:p03-w03-p03-047-q0-not-re}.

(4) For \(n\ge1\) let \(\delta_n\) say that there are at least \(n\)
elements.  Every finite subset of
\(\{\exists X\,\forall x\,Xx\}\cup\{\delta_n:n\ge1\}\) holds in a finite
structure of power the largest \(n\) mentioned (or \(1\)), while a model of
the whole set would be finite with at least \(n\) elements for every \(n\).
For validity: by part~3 and the argument of part~2 of
\cref{thm:p03-w03-p03-047-q0-not-re}, \(\mathcal N\models\sigma\) iff
\((\chi_{\mathcal N})^w\to\sigma\) is valid, for every first-order
\(\Lang_A\)-sentence \(\sigma\).  If the valid sentences were recursively
enumerable, the computable map \(\sigma\mapsto(\chi_{\mathcal N})^w\to\sigma\)
would make \(\mathrm{TA}\) recursively enumerable, exactly as in part~3 of
that proof, contradicting clause~3 of
\cref{thm:p02-w02-p02-040-no-effective-calculus}.

(5) For a formula \(\varphi\) whose free relation variables are among
\(X_1,\ldots,X_m\), \(X_j\) of arity \(n_j\), and for each
\(\bar k=(k_1,\ldots,k_m)\in\omega^m\), choose individual variables
\(\bar y^j_i=(y^j_{i,1},\ldots,y^j_{i,n_j})\) (\(1\le i\le k_j\)) not
occurring in \(\varphi\), and define an \(\mathscr L_{\omega_1\omega}\)-formula
\(\varphi_{\bar k}\) by recursion on \(\varphi\): an atomic formula without
relation variables is unchanged; \(X_j\bar t\) becomes
\(\bigvee_{i=1}^{k_j}\bigwedge_{l=1}^{n_j}t_l=y^j_{i,l}\), read as
\(\exists x\,\neg x=x\) (false, domains being nonempty) when \(k_j=0\);
\(\varphi_{\bar k}\) commutes with \(\neg,\land,\lor\) and with \(\exists x\),
\(\forall x\); and for a new \(n\)-ary variable \(X\),
\[
 (\exists X\,\psi)_{\bar k}=\bigvee_{k\in\omega}\exists\bar y_1\ldots\exists\bar y_k\,\psi_{\bar k,k},
 \qquad
 (\forall X\,\psi)_{\bar k}=\bigwedge_{k\in\omega}\forall\bar y_1\ldots\forall\bar y_k\,\psi_{\bar k,k},
\]
with fresh \(\bar y_i\) of length \(n\).  Each such formula has countable
conjunctions and disjunctions and finite quantifier strings, so lies in
\(\mathscr L_{\omega_1\omega}\).

Claim: for every structure \(\mathfrak A\) and assignment \(s\),
\(\mathfrak A,s[X_j\mapsto V_j]\models\varphi\) iff
\(\mathfrak A,s\models\varphi_{\bar k}\), where
\(V_j=\{s(\bar y^j_1),\ldots,s(\bar y^j_{k_j})\}\).  Induction on \(\varphi\).
For \(X_j\bar t\): the value of \(\bar t\) lies in \(V_j\) iff it equals some
\(s(\bar y^j_i)\).  The cases of connectives and individual quantifiers are
immediate from the induction hypothesis, since the variables \(\bar y\) are
not quantified in \(\varphi\).  For \(\exists X\,\psi\): the left side holds
iff some finite \(n\)-ary \(V\) has \(\mathfrak A,s[\bar X\mapsto\bar V,X\mapsto V]\models\psi\).
Every finite \(V\) is \(\{\bar b_1,\ldots,\bar b_k\}\) for \(k=|V|\) and some
tuples, and every such set is finite; so the left side holds iff for some
\(k\) and tuples \(\bar b_1,\ldots,\bar b_k\) the formula \(\psi\) holds with
\(X\mapsto\{\bar b_1,\ldots,\bar b_k\}\), iff (induction hypothesis, applied
to \(s\) modified on \(\bar y_1,\ldots,\bar y_k\)) the disjunct for \(k\)
holds.  The case \(\forall X\) is dual.

For \(\varphi\) without free relation variables put
\(\varphi^\ast=\varphi_{()}\), the case \(m=0\).  If \(\varphi\) has a model,
so does the \(\mathscr L_{\omega_1\omega}\)-sentence \(\varphi^\ast\), which
by part~1 of \cref{thm:p02-w02-p02-005-bw-extended-ls} has a model of power
at most \(\aleph_0\); by the claim that model satisfies \(\varphi\).
\end{proof}
\end{annexitem}

\begin{annexitem}{annex:p03-w03-p03-047-q1-ls-consequence}{thm:p03-w03-p03-047-q1-ls-consequence}{Downward L\"owenheim--Skolem to \(\aleph_1\) and consequence for \(\mathscr L(Q_1)\)}
\begin{proof}
(1) Let \(\mathfrak A\models\Phi\) and let \(\tau_0\) be the set of symbols
occurring in members of \(\Phi\); \(|\tau_0|\le\aleph_1\) because each
sentence has finitely many symbols.  By the reduct property work with
\(\mathfrak A_0=\mathfrak A\restriction\tau_0\).  There are at most
\(\aleph_1\) formulas of \(\mathscr L(Q_1)[\tau_0]\).  Fix a well-ordering of
\(A\) to make choices.  Build \(B_\alpha\subseteq A\), \(\alpha<\omega_1\),
increasing, each of power at most \(\aleph_1\).  \(B_0\) consists of one
element and the interpretations of the constants of \(\tau_0\).  Given
\(B_\alpha\), let \(B'_\alpha\) be \(B_\alpha\) together with, for each
formula \(\varphi(x,\bar y)\) and each \(\bar b\) from \(B_\alpha\):
(a) if \(\mathfrak A_0\models\exists x\,\varphi[\bar b]\), the least
\(a\) with \(\mathfrak A_0\models\varphi[a,\bar b]\); (b) if
\(\mathfrak A_0\models Q_1x\,\varphi[\bar b]\) and some \(a\notin B_\alpha\)
has \(\mathfrak A_0\models\varphi[a,\bar b]\), the least such \(a\).  Let
\(B_{\alpha+1}\) be the closure of \(B'_\alpha\) under the functions of
\(\tau_0\).  At limits take unions, and let \(B=\bigcup_{\alpha<\omega_1}B_\alpha\).
There are at most \(\aleph_1\cdot\aleph_1=\aleph_1\) pairs of a formula and a
finite tuple from a set of power \(\aleph_1\), and closing a set of power at
most \(\aleph_1\) under at most \(\aleph_1\) finitary functions gives a set of
power at most \(\aleph_1\) (the closure is the union over \(n\in\omega\) of the
\(n\)-step closures); by induction every \(B_\alpha\) has power at most
\(\aleph_1\), and so does \(B\).  A finite tuple from \(B\) lies in one
\(B_\alpha\), so \(B\) is closed under the functions; it contains the
constants and is nonempty, so it carries a substructure \(\mathfrak B\) of
\(\mathfrak A_0\).

Claim: for every formula \(\varphi(\bar y)\) and \(\bar b\) from \(B\),
\(\mathfrak B\models\varphi[\bar b]\) iff \(\mathfrak A_0\models\varphi[\bar b]\).
Atomic formulas are preserved by substructures; the Boolean steps are
immediate.  For \(\exists x\,\varphi\): if \(\mathfrak B\models\varphi[a,\bar b]\)
then \(\mathfrak A_0\models\varphi[a,\bar b]\) by induction; if
\(\mathfrak A_0\models\exists x\,\varphi[\bar b]\), with \(\bar b\) in
\(B_\alpha\), step (a) put a witness into \(B_{\alpha+1}\subseteq B\), and
the induction hypothesis transfers it.  For \(Q_1x\,\varphi\): let
\(W=\{a\in A:\mathfrak A_0\models\varphi[a,\bar b]\}\); by induction the
corresponding set in \(\mathfrak B\) is \(W\cap B\).  If \(W\) is countable,
so is \(W\cap B\), and both sides fail.  If \(W\) is uncountable, let
\(\bar b\) lie in \(B_{\alpha_0}\).  If \(W\subseteq B_\beta\) for some
\(\beta\), then \(W\cap B=W\) is uncountable.  Otherwise, for each
\(\beta\ge\alpha_0\), step (b) at stage \(\beta\) adds some
\(w_\beta\in W\cap B_{\beta+1}\setminus B_\beta\).  For
\(\beta<\gamma\), \(w_\beta\in B_{\beta+1}\subseteq B_\gamma\) while
\(w_\gamma\notin B_\gamma\), so the \(w_\beta\) are distinct, and
\(W\cap B\) has power \(\aleph_1\).  In both cases both sides hold.

So \(\mathfrak B\models\Phi\) and \(|B|\le\aleph_1\).  Expand \(\mathfrak B\)
to \(\tau_{\mathfrak A}\) by interpreting the remaining symbols arbitrarily
on the nonempty domain; by the reduct property the expansion is a model of
\(\Phi\) of power at most \(\aleph_1\).

(2) For a sentence \(\varphi\), \(\Sigma\models\varphi\) iff
\(\Sigma\cup\{\neg\varphi\}\) has no standard model, iff (by
\cref{thm:p02-w02c-p02-003-keisler-completeness}) it is inconsistent, iff for
some sentence \(\sigma\) there are derivations of \(\sigma\) and of
\(\neg\sigma\) from \(\Sigma\cup\{\neg\varphi\}\) in the sense of
\cref{def:p02-w02c-p02-003-kq-calculus}.  Since \(\Sigma\) is recursively
enumerable and the axioms form a decidable set, the set of triples
\((\varphi,d_1,d_2)\) with \(d_1,d_2\) such derivations, ending in some
\(\sigma\) and in \(\neg\sigma\), is recursively enumerable, and its
projection to \(\varphi\) is the set of consequences of \(\Sigma\).
\end{proof}
\end{annexitem}

\begin{annexitem}{annex:p03-w03-p03-047-q1-interpolants}{thm:p03-w03-p03-047-q1-interpolants}{Interpolants in \(\mathscr L(Q_1^2)\) and \(\mathscr L(\mathrm{pos})\)}
\begin{proof}
\emph{The entailment.}  In a model of \(\varphi_0\land\varphi_1\), \(E\) is an
equivalence relation and \(R\), \(S\) are sets of representatives: each class
contains exactly one element of each.  The map sending \(r\in R\) to the
element of \(S\) in its class is a bijection of \(R\) onto \(S\), so
\(R\) countable and \(S\) uncountable is impossible; hence
\(\varphi_0\models\neg\varphi_1\).  The same argument shows: if \(E\) is an
equivalence relation with a set \(P\) of representatives, \(P\) is in
bijection with the set of \(E\)-classes.

(1) If \(\psi^{\mathfrak A}\) is uncountable, \(M=\psi^{\mathfrak A}\) gives
\(\psi(b_0)\land\psi(b_1)\) for all \((b_0,b_1)\in M^2\).  If an uncountable
\(M\) has this property, then for \(b\in M\) the pair \((b,b)\) gives
\(\psi(b)\), so \(M\subseteq\psi^{\mathfrak A}\) and \(\psi^{\mathfrak A}\)
is uncountable.  Replacing in \(\varphi(E)\) each \(Q_1y\,Exy\) by
\(Q_1^2yz\,(Exy\land Exz)\), with \(z\) fresh, gives an equivalent
sentence of \(\mathscr L_{\omega\omega}(Q_1^2)[\{E\}]\).

Let \(E\) be an equivalence relation on \(A\).  If an uncountable \(M\) has
\(b_0=b_1\lor\neg Eb_0b_1\) for all pairs from \(M\), its elements lie in
distinct classes, so there are uncountably many classes.  If there are
uncountably many classes, choose one element from each (axiom of choice);
the set \(M\) of chosen elements is uncountable and any two distinct members
are inequivalent.  So \(Q_1^2xy\,(x=y\lor\neg Exy)\) holds iff \(E\) has
uncountably many classes.  Let \(\psi=\varphi(E)\land\neg Q_1^2xy\,(x=y\lor\neg Exy)\).
In a model of \(\varphi_0\), \(R\) is a countable set of representatives, so
there are countably many classes and \(\psi\) holds.  In a model of \(\psi\)
there are countably many classes, so no set of representatives is
uncountable and \(\varphi_1\) fails.  Thus \(\varphi_0\models\psi\) and
\(\psi\models\neg\varphi_1\), with \(\psi\) in the common vocabulary
\(\{E\}\).

(2) Read \(Q_1y\,Exy\) as \(\neg\exists X\,\forall y\,(Xy\lor\neg Exy)\), with
\(X\) over countable sets: \(\exists X\,\forall y\,(Xy\lor\neg Exy)\) holds
iff some countable set covers the class of \(x\), iff that class is
countable.  The single occurrence of \(X\) in the matrix is in the scope of
no negation, so the quantification is admitted in
\(\mathscr L_{\omega\omega}(\mathrm{pos})\); the remaining clauses of
\(\varphi(E)\) are first-order.  In \(\exists X\,\forall y\,\exists x\,(Xx\land Exy)\)
the variable \(X\) again occurs under no negation.  For an equivalence
relation \(E\), this sentence holds iff some countable \(X\) meets every
class.  If so, choosing for each class an element of \(X\) in it is an
injection of the set of classes into \(X\), so there are countably many
classes; if there are countably many classes, a choice of representatives is
such an \(X\).  So the sentence \(\psi'=\varphi(E)\land\exists X\,\forall y\,\exists x\,(Xx\land Exy)\)
says, given \(\varphi(E)\), that there are countably many classes, and the
argument of part~(1) gives \(\varphi_0\models\psi'\models\neg\varphi_1\).
\end{proof}
\end{annexitem}

\begin{annexitem}{annex:p03-w03-p03-047-set-quantifier-failures}{thm:p03-w03-p03-047-set-quantifier-failures}{Characterizing \((\omega_1,<)\) and \((\omega,<)\) with set quantifiers}
\begin{proof}
(1) A linear ordering \((A,<)\) is \(\aleph_1\)-like iff \(A\) is
uncountable and each set \(\{z:z<y\}\) is countable.  \((\omega_1,<)\) is
\(\aleph_1\)-like, since every countable ordinal has countably many
predecessors, and it is a well-ordering, so every nonempty subset has a least
element.  Conversely let \((A,<)\) be \(\aleph_1\)-like with least elements
for uncountable subsets, and let \(Y\subseteq A\) be nonempty.  If \(Y\) is
uncountable it has a least element.  If \(Y\) is countable, let
\(U=\{z\in A:y<z\text{ for all }y\in Y\}\).  Its complement is
\(\bigcup_{y\in Y}\{z:z\le y\}\), a countable union of countable sets, hence
countable (axiom of choice); so \(U\) is uncountable, and \(Y\cup U\) is
uncountable and has a least element \(m\).  If \(m\in U\), any \(y\in Y\)
satisfies \(y<m\), contradicting minimality; so \(m\in Y\), and \(m\) is
least in \(Y\).  Thus \((A,<)\) is a well-ordering, isomorphic to an ordinal
\(\alpha\).  As \(A\) is uncountable, \(\alpha\ge\omega_1\); if
\(\alpha>\omega_1\), the element at position \(\omega_1\) would have
uncountably many predecessors.  So \(\alpha=\omega_1\).

The sentence: ``\(<\) is a strict linear order''
\(\land\,Q_1x\,x=x\land\forall y\,\neg Q_1x\,x<y\,\land\,
\neg\exists\text{ uncountable }X\,\neg\exists x\,(Xx\land\forall y\,(\neg Xy\lor\neg y<x))\),
with \(Q_1\) expressed by an uncountable-set quantifier as above; by what was
just shown its models are the copies of \((\omega_1,<)\).  Add constants
\(c_n\), \(n\in\omega\), and the sentences \(c_{n+1}<c_n\).  A finite subset
mentions \(c_0,\ldots,c_N\) only; in \((\omega_1,<)\) interpret \(c_n\) as
\(N-n\) for \(n\le N\) and every other constant as \(0\).  The whole set has no
model, since a model would be a well-ordering in which \(\{c_n\}\) has no
least element.  This countable set witnesses the failure of
\(\aleph_0\)-compactness.

(2) Let the sentence say: \(<\) is a strict linear order;
\(\exists\) countable \(X\,\forall x\,Xx\); for every countable \(X\), if
\(\exists x\,Xx\) then \(X\) has a least element; \(\forall x\,\exists y\,x<y\);
and every element with a predecessor has an immediate predecessor.  In a
model the domain is countable, so every subset is countable and has a least
element if nonempty: the order is a well-ordering of type some \(\alpha\),
which is nonzero, has no largest element, and has every nonzero member a
successor.  As in the proof for \(\Omega\) in
\cref{thm:p03-w03-p03-047-sol-ch}, \(\alpha=\omega\).  The copies of
\((\omega,<)\) satisfy every clause.  Add a constant \(c\) and the sentences
\(\delta_n\) of the proof of \cref{thm:p03-w03-p03-047-q1-mostowski-facts}
saying that \(c\) has at least \(n\) predecessors; every finite subset holds
in \((\omega,<)\) with \(c\) interpreted as a large enough number, and the whole
set has no model, since in a copy of \((\omega,<)\) every element has
finitely many predecessors.
\end{proof}
\end{annexitem}

\begin{annexitem}{annex:p03-w03-p03-047-pos-monotone}{thm:p03-w03-p03-047-pos-monotone}{Monotonicity in positive set variables}
\begin{proof}
Prove, by induction on \(\varphi\) and simultaneously for all set variables
\(X\) free in \(\varphi\): (a) if \(X\) has no negative occurrence in
\(\varphi\), then \(\varphi\) is upward monotone in \(X\) over countable
sets; (b) if \(X\) has no positive occurrence, \(\varphi\) is downward
monotone in \(X\).  The other free variables keep fixed values throughout.

If \(\varphi\) is \(Xt\) for a term \(t\), the
occurrence is positive, and if the value of \(t\) lies in \(M\) it lies in
\(M'\supseteq M\): (a) holds and (b) is vacuous.  Atomic formulas without
\(X\) do not change.  If \(\varphi=\neg\psi\), an occurrence of \(X\) is
positive in \(\varphi\) iff it is negative in \(\psi\); under (a) for
\(\varphi\), \(\psi\) is downward monotone by (b) for \(\psi\), so
\(\neg\psi\) is upward monotone; (b) is symmetric.  For \(\psi_0\land\psi_1\)
and \(\psi_0\lor\psi_1\), each occurrence keeps its polarity, so both
components are monotone in the same direction, and so are their conjunction
and disjunction.  For \(\exists x\,\psi\) and \(\forall x\,\psi\): for each
value of \(x\), \(\psi\) is monotone in the given direction, and the
existential and universal closures of a family of monotone conditions are
monotone.  For \(\exists Y\,\psi\) with \(Y\ne X\): polarities of \(X\) are
unchanged; if \(\mathfrak A\models\exists Y\,\psi[M,\ldots]\) under (a), some
countable \(N\) has \(\mathfrak A\models\psi[M,N,\ldots]\), hence
\(\mathfrak A\models\psi[M',N,\ldots]\) and
\(\mathfrak A\models\exists Y\,\psi[M',\ldots]\); (b) is symmetric.  If
\(X\) is bound by the outermost quantifier it is not free, and there is
nothing to prove.  The statement of the theorem is (a) and (b).
\end{proof}
\end{annexitem}

\begin{annexitem}{annex:p03-w03-p03-047-pc-basic}{thm:p03-w03-p03-047-pc-basic}{PC classes inside RPC classes; infinity}
\begin{proof}
(1) Let \(K\) be PC in \(\mathscr L\), witnessed by \(\tau'\supseteq\tau\)
and \(\varphi\in\mathscr L[\tau']\) with
\(K=\{\mathfrak A\restriction\tau:\mathfrak A\models\varphi\}\).  By the
renaming property we may assume that a unary symbol \(U\) does not occur in
\(\tau'\); put \(\tau''=\tau'\cup\{U\}\).  By hypothesis there is
\(\upsilon\in\mathscr L[\tau'']\) whose models are the
\(\tau''\)-structures with \(U^{\mathfrak A}=A\), and the conjunction
property supplies \(\chi\in\mathscr L[\tau'']\) equivalent to
\(\varphi\land\upsilon\) (with \(\varphi\in\mathscr L[\tau'']\) by
property~(i) of \cref{def:p03-w03-p03-047-ebbinghaus-logic}).  Let
\(K''=\operatorname{Mod}_{\mathscr L}(\chi)\).  For \(\mathfrak A\in K''\),
\(U^{\mathfrak A}=A\) is \(\tau\)-closed and
\((\mathfrak A\restriction\tau)|U^{\mathfrak A}=\mathfrak A\restriction\tau\),
and \(\mathfrak A\restriction\tau'\models\varphi\) by the reduct property, so
\(\mathfrak A\restriction\tau\in K\).  Conversely, if
\(\mathfrak B\in K\), say \(\mathfrak B=\mathfrak A\restriction\tau\) with
\(\mathfrak A\models\varphi\), expand \(\mathfrak A\) by \(U=A\); the
expansion lies in \(K''\) and yields \(\mathfrak B\).  So \(K\) is RPC.

(2) A set \(A\) carries an injective non-surjective \(f:A\to A\) iff \(A\) is
infinite: a finite set has no such map (an injection of a finite set into
itself is surjective), and an infinite \(A\) contains distinct
\(a_0,a_1,\ldots\) (in \(\mathsf{ZFC}\)), and the map with
\(a_n\mapsto a_{n+1}\) and identity elsewhere is injective and misses
\(a_0\).  The displayed sentence says exactly that \(f\) is injective and not
surjective, so its reducts to the empty vocabulary are the infinite sets.
Suppose some first-order sentence \(\psi\) of the empty vocabulary held
exactly in the infinite sets.  Let \(\lambda_n\) be the first-order sentence
``there are at least \(n\) elements''.  Every finite subset of
\(\{\neg\psi\}\cup\{\lambda_n:n\ge1\}\) has a model, namely a finite set
of size at least the largest \(n\) involved, which satisfies \(\neg\psi\).
By the compactness theorem (\cref{thm:p02-w02-p02-005-bw-compactness}) the
whole set has a model, which is infinite by the \(\lambda_n\) and satisfies
\(\neg\psi\), a contradiction.
\end{proof}
\end{annexitem}

\begin{annexitem}{annex:p03-w03-p03-047-strength-basic}{thm:p03-w03-p03-047-strength-basic}{Comparison relations}
\begin{proof}
(1) Every class EC in \(\mathscr L\) is EC in \(\mathscr L\), and if every
EC class of \(\mathscr L\) is EC in \(\mathscr L'\) and every EC class of
\(\mathscr L'\) is EC in \(\mathscr L''\), then every EC class of
\(\mathscr L\) is EC in \(\mathscr L''\).  A reflexive transitive relation
makes its symmetric part an equivalence relation and induces a partial
ordering of the equivalence classes, well defined because \(\le\) is
transitive.  The same argument applies word for word with PC or RPC in place
of EC.

(2) Let \(\mathscr L\le\mathscr L^*\), let
\(\mathfrak A,\mathfrak B\in\operatorname{Str}[\tau]\) with
\(\mathfrak A\equiv_{\mathscr L^*}\mathfrak B\), and let
\(\varphi\in\mathscr L[\tau]\).  The class
\(\operatorname{Mod}_{\mathscr L}(\varphi)\) of \(\tau\)-structures is EC in
\(\mathscr L^*\), so some \(\varphi^*\in\mathscr L^*[\tau]\) has
\(\operatorname{Mod}_{\mathscr L^*}(\varphi^*)=\operatorname{Mod}_{\mathscr L}(\varphi)\).
Then \(\mathfrak A\models\varphi\) iff \(\mathfrak A\models\varphi^*\) iff
\(\mathfrak B\models\varphi^*\) iff \(\mathfrak B\models\varphi\).  Hence
\(\mathfrak A\equiv_{\mathscr L}\mathfrak B\).
\end{proof}
\end{annexitem}

\begin{annexitem}{annex:p03-w03-p03-047-rpc-descent}{thm:p03-w03-p03-047-rpc-descent}{Descent of compactness and L\"owenheim--Skolem}
\begin{proof}
For \(\varphi\in\mathscr L[\tau]\), the class
\(\operatorname{Mod}_{\mathscr L}(\varphi)\) is EC, hence RPC, in
\(\mathscr L\), hence RPC in \(\mathscr L^*\) in the many-sorted reading:
there are \(\tau_\varphi\supseteq\tau\) and
\(\varphi^*\in\mathscr L^*[\tau_\varphi]\) with
\[
 \operatorname{Mod}_{\mathscr L}(\varphi)
 =\{\mathfrak C\restriction\tau:\mathfrak C\models_{\mathscr L^*}\varphi^*\}. \tag{\(\dagger\)}
\]

(1) Let \(\Phi\subseteq\mathscr L[\tau]\) have power at most \(\kappa\) with
every finite subset satisfiable.  For each \(\varphi\in\Phi\) choose
\(\tau_\varphi\) and \(\varphi^*\) as in (\(\dagger\)).  By the renaming
property of \(\mathscr L^*\), applied to a renaming that fixes \(\tau\) and
sends \(\tau_\varphi\setminus\tau\) to fresh symbols (sort symbols to sort
symbols), we may assume that the sets \(\tau_\varphi\setminus\tau\) for
distinct \(\varphi\in\Phi\) are pairwise disjoint; a renaming fixing
\(\tau\) does not change \(\tau\)-reducts, so (\(\dagger\)) persists.  Let
\(\tau^*=\bigcup_{\varphi\in\Phi}\tau_\varphi\) and
\(\Phi^*=\{\varphi^*:\varphi\in\Phi\}\subseteq\mathscr L^*[\tau^*]\), of
power at most \(\kappa\).  Let \(\Phi_0^*\subseteq\Phi^*\) be finite, coming
from a finite \(\Phi_0\subseteq\Phi\), and let \(\mathfrak A\models\Phi_0\).
For each \(\varphi\in\Phi_0\), (\(\dagger\)) gives a
\(\tau_\varphi\)-structure \(\mathfrak C_\varphi\models\varphi^*\) with
\(\mathfrak C_\varphi\restriction\tau=\mathfrak A\).  Because the new parts
\(\tau_\varphi\setminus\tau\) are pairwise disjoint, there is a
\(\tau^*\)-structure \(\mathfrak D\) with
\(\mathfrak D\restriction\tau_\varphi=\mathfrak C_\varphi\) for
\(\varphi\in\Phi_0\): it agrees with \(\mathfrak A\) on \(\tau\), with
\(\mathfrak C_\varphi\) on the symbols and new sorts of
\(\tau_\varphi\setminus\tau\), and interprets the remaining symbols
arbitrarily over nonempty new domains.  By the reduct property
\(\mathfrak D\models\Phi_0^*\).  So every finite subset of \(\Phi^*\) is
satisfiable, and \(\kappa\)-compactness of \(\mathscr L^*\) gives
\(\mathfrak E\models\Phi^*\).  For \(\varphi\in\Phi\),
\(\mathfrak E\restriction\tau_\varphi\models\varphi^*\) by the reduct
property, so \(\mathfrak E\restriction\tau\models\varphi\) by (\(\dagger\)).
Thus \(\Phi\) is satisfiable.  If \(\mathscr L^*\) is compact it is
\(\kappa\)-compact for every infinite \(\kappa\), hence so is
\(\mathscr L\).

(2) Let \(\varphi\in\mathscr L[\tau]\) be satisfiable.  By (\(\dagger\)) the
sentence \(\varphi^*\) is satisfiable, so it has a model \(\mathfrak C\) of
power at most \(\kappa\).  Its reduct \(\mathfrak C\restriction\tau\) is a
model of \(\varphi\) by (\(\dagger\)), and its power is a sum over the sorts
of \(\tau\) of the same cardinals that occur in the sum for \(\mathfrak C\),
so it is at most \(\kappa\).
\end{proof}
\end{annexitem}

\begin{annexitem}{annex:p03-w03-p03-047-robinson-interpolation}{thm:p03-w03-p03-047-robinson-interpolation}{Robinson property and interpolation}
\begin{proof}
Since \(\mathscr L\) is small, every class of sentences formed below is a
set, so compactness applies to it.  For a finite set \(\Psi\) of sentences
of one vocabulary, \(\bigwedge\Psi\) denotes a sentence given by the
conjunction property with the same models as \(\Psi\) (for
\(\Psi=\emptyset\), the valid sentence \(\exists c\,(c=c)\) supplied by the
atom and particularization properties).  By the reduct remark at the head of this annex, consequence
between sentences of different vocabularies may be evaluated over any
vocabulary containing both.

\emph{Robinson implies interpolation.}  Let \(\varphi_i\in\mathscr L[\tau_i]\)
(\(i=0,1\)) with \(\varphi_0\models\varphi_1\), and \(\tau=\tau_0\cap\tau_1\).
Put \(\Phi'=\{\varphi\in\mathscr L[\tau]:\varphi_0\models\varphi\}\).  We
claim \(\Phi'\models\varphi_1\).  If not, some
\((\tau_1\cup\tau)\)-structure \(\mathfrak C\) satisfies
\(\Phi'\cup\{\neg\varphi_1\}\); let \(\mathfrak B=\mathfrak C\restriction\tau\)
and \(\Phi=\operatorname{Th}_{\mathscr L}(\mathfrak B)\), the set of
\(\varphi\in\mathscr L[\tau]\) true in \(\mathfrak B\).  \(\Phi\) is
complete: a \(\tau\)-model of \(\Phi\) satisfies, for each
\(\varphi\in\mathscr L[\tau]\), either \(\varphi\) or its negation as
\(\mathfrak B\) does.  \(\Phi\cup\{\neg\varphi_1\}\) has the model
\(\mathfrak C\).  \(\Phi\cup\{\varphi_0\}\) has a model: otherwise, by
compactness, a finite \(\Psi\subseteq\Phi\) has \(\Psi\cup\{\varphi_0\}\)
unsatisfiable, so \(\varphi_0\models\neg\bigwedge\Psi\) and
\(\neg\bigwedge\Psi\in\Phi'\); then \(\mathfrak B\models\neg\bigwedge\Psi\),
contradicting \(\Psi\subseteq\Phi\).  The Robinson property, with
\(\Phi_0=\{\varphi_0\}\) and \(\Phi_1=\{\neg\varphi_1\}\), gives a model of
\(\{\varphi_0,\neg\varphi_1\}\), contradicting
\(\varphi_0\models\varphi_1\).  So \(\Phi'\models\varphi_1\); by
compactness some finite \(\Phi''\subseteq\Phi'\) has
\(\Phi''\models\varphi_1\).  Then \(\psi=\bigwedge\Phi''\in\mathscr L[\tau]\)
satisfies \(\varphi_0\models\psi\) (each member of \(\Phi''\) follows from
\(\varphi_0\)) and \(\psi\models\varphi_1\).

\emph{Interpolation implies Robinson.}  Let \(\Phi,\Phi_0,\Phi_1\) be as in
\cref{def:p03-w03-p03-047-robinson-property}, with \(\Phi\) complete and
\(\Phi\cup\Phi_i\) satisfiable.  Replacing \(\Phi_0\) by \(\Phi\cup\Phi_0\)
(\(\tau\subseteq\tau_0\)) we may assume \(\Phi\subseteq\Phi_0\).  By
compactness it suffices that for every finite \(\Psi_i\subseteq\Phi_i\) the
set \(\{\varphi_0,\varphi_1\}\) with \(\varphi_i=\bigwedge\Psi_i\) has a
model.  Suppose not.  Then \(\varphi_0\models\neg\varphi_1\), and
interpolation gives \(\psi\in\mathscr L[\tau]\) with \(\varphi_0\models\psi\)
and \(\psi\models\neg\varphi_1\).  A model of \(\Phi_0\) is a model of
\(\Phi\cup\{\varphi_0\}\) and so of \(\psi\); its \(\tau\)-reduct is a model
of \(\Phi\) satisfying \(\psi\), and since all \(\tau\)-models of \(\Phi\)
are \(\mathscr L\)-equivalent, \(\Phi\models\psi\).  Then
\(\Phi\models\neg\varphi_1\), so \(\Phi\cup\Phi_1\) has no model, a
contradiction.
\end{proof}
\end{annexitem}

\begin{annexitem}{annex:p03-w03-p03-047-infinite-sets}{thm:p03-w03-p03-047-infinite-sets}{Equivalence of infinite sets; the two counterexamples}
\begin{proof}
(1) Two tuples \(\bar a=(a_1,\dots,a_k)\) of \(A\) and
\(\bar b=(b_1,\dots,b_k)\) of \(B\) have the \emph{same equality type} if
\(a_i=a_j\) iff \(b_i=b_j\) for all \(i,j\).  Let \(A\) and \(B\) be infinite
sets.  We show by induction on formulas \(\varphi(x_1,\dots,x_k)\) of
\(\mathscr L_{\infty\omega}\) of the empty vocabulary that
\(A\models\varphi[\bar a]\) iff \(B\models\varphi[\bar b]\) whenever
\(\bar a\) and \(\bar b\) have the same equality type.  The atomic formulas
are equations \(x_i=x_j\), for which this is the definition.  Negation and
conjunctions or disjunctions of any set of formulas pass through the
induction hypothesis.  Let \(\varphi=\exists y\,\chi(\bar x,y)\) and
\(A\models\chi[\bar a,a']\).  If \(a'=a_i\) for some \(i\), put \(b'=b_i\);
otherwise choose \(b'\in B\setminus\{b_1,\dots,b_k\}\), possible because
\(B\) is infinite.  Then \((\bar a,a')\) and \((\bar b,b')\) have the same
equality type, so \(B\models\chi[\bar b,b']\) by the induction hypothesis.
The converse is symmetric.  A homogeneous string of finitely many
quantifiers is a finite iteration of this case.  For sentences
(\(k=0\)) this gives part~(1).

(2) Let \(\{c_n:n<\omega\}\) and \(\{d_\alpha:\alpha<\omega_1\}\) be sets of
distinct new constants.  \(\mathfrak K_{\aleph_0}\) is the class of
reducts to the empty vocabulary of the models of
\(\forall x\bigvee_{n<\omega}x=c_n\): a model of it has power at most
\(\aleph_0\), and a nonempty set of power at most \(\aleph_0\) is the image
of a map \(n\mapsto c_n\) from \(\omega\).  \(\mathfrak K^{\aleph_1}\) is
the class of reducts of the models of
\(\bigwedge_{\alpha<\beta<\omega_1}\neg d_\alpha=d_\beta\), a conjunction of
\(\aleph_1\) formulas, hence a sentence of \(\mathscr L_{\omega_2\omega}\);
a set of power at least \(\aleph_1\) has \(\aleph_1\) distinct elements to
interpret the \(d_\alpha\).  The first sentence is in
\(\mathscr L_{\omega_1\omega}\subseteq\mathscr L_{\omega_2\omega}\), so both
classes are PC in \(\mathscr L_{\omega_2\omega}\), and they are disjoint.
If a class \(\operatorname{Mod}(\psi)\) with \(\psi\) an
\(\mathscr L_{\infty\omega}\)-sentence of the empty vocabulary contained
\(\mathfrak K_{\aleph_0}\) and were disjoint from
\(\mathfrak K^{\aleph_1}\), then \(\omega\models\psi\) and
\(\omega_1\not\models\psi\), contradicting part~(1).

(3) Use the constants of part~(2) and take \(\tau_0=\{c_n:n<\omega\}\),
\(\tau_1=\{d_\alpha:\alpha<\omega_1\}\), so \(\tau=\tau_0\cap\tau_1=\emptyset\).
Let \(\Phi=\{\iota\}\), where \(\iota\) is the
\(\mathscr L_{\omega_1\omega}\)-sentence
\(\bigwedge_{n\ge1}\exists x_1\cdots\exists x_n\bigwedge_{1\le i<j\le n}\neg x_i=x_j\),
true exactly in the infinite sets.  \(\Phi\) is complete by part~(1),
since \(\mathscr L_{\omega_1\omega}\subseteq\mathscr L_{\infty\omega}\).
Let \(\Phi_0=\{\forall x\bigvee_nx=c_n\}\) and
\(\Phi_1=\{\neg d_\alpha=d_\beta:\alpha<\beta<\omega_1\}\), a set of
\(\mathscr L_{\omega_1\omega}[\tau_1]\)-sentences.  \(\Phi\cup\Phi_0\) has
the model \(\omega\) with \(c_n\) read as \(n\); \(\Phi\cup\Phi_1\) has the
model \(\omega_1\) with \(d_\alpha\) read as \(\alpha\).  A model of
\(\Phi\cup\Phi_0\cup\Phi_1\) would have power at most \(\aleph_0\) by
\(\Phi_0\) and at least \(\aleph_1\) by \(\Phi_1\).  So the Robinson
property fails for \(\mathscr L_{\omega_1\omega}\).
\end{proof}
\end{annexitem}
